\documentclass{llncs}
\usepackage{hyperref}
% \usepackage[fleqn]{amsmath}
\usepackage{amsmath}
\usepackage{amsfonts}
\usepackage{float}
%\usepackage{amsthm}
\usepackage{graphicx}
\usepackage{booktabs}   % For professional-quality tables
\usepackage[T1]{fontenc} %(or \usepackage{t1enc}) 

\usepackage{tikz}
\usepackage{enumitem}
\usepackage{multirow} 

\usepackage[cache=false]{minted}
%\usepackage{minted}


\setminted{
  % Line numbers not flowing out of the margin
  numbersep=5pt,
  xleftmargin=12pt,
  %
  % Better listing breaking
  breakafter=-/\{,
  breakbefore=\\
  %
  % Alternative: Rely on pygment's tokenizer. Does not work well with LaTeX and comments
  % breakbytoken=true,
  % breakbytokenanywhere=true
}

\usepackage{listings}  % Include the listings package for code formatting


% Define SMT-LIB language for syntax highlighting
\lstdefinelanguage{SMT-LIB}{
  morekeywords={
    declare-const, assert, check-sat, define-fun, forall, or, and, not, 
    exists, if, then, else, let, in, re, str, seq, Int, String, 
    exists, true, false, "=",
  },
  sensitive=true,
  morecomment=[l]{;}, % Single-line comments
}

\lstset{ %
  language=java,               % Use 'java' for basic highlighting
  basicstyle=\ttfamily\small,   % Use a small monospaced font
  numbers=left,                % Show line numbers on the left
  numberstyle=\tiny\color{gray},% Line number style
  stepnumber=1,                % Number every line
  numbersep=5pt,               % Distance between line numbers and code
  backgroundcolor=\color{white},% Background color for the code
  showspaces=false,            % Do not highlight spaces
  showstringspaces=false,      % Do not highlight spaces in strings
  showtabs=false,              % Do not highlight tabs
  frame=single,                % Frame around the code block
  rulecolor=\color{black},     % Color of the rule/frame
  tabsize=2,                   % Tab size
  captionpos=b,                % Position of the caption (b for bottom)
  breaklines=true,             % Break long lines
  breakatwhitespace=false,     % Break at whitespace
  escapeinside={\%*}{*)},      % Allows you to escape to LaTeX
}

%%%%%%%%%%%%%%%%%%%%%%%%%%%%%%%%%%%%%%%%%%%%%%%%%%%%%%%%%%%%%%%%%%%%%%%%%%%%%%%%%%%%%%%

\newcommand{\cerl}{\mathcal{R}}
\newcommand*{\aut}{\mathcal{A}}
\newcommand*{\NFA}{\mathcal{A}}
\newcommand*{\autb}{\mathcal{B}}
\newcommand*{\transducer}{\mathfrak{T}}
\newcommand*{\Int}{\mathbb{Z}}
\newcommand*{\myvec}[1]{\overrightarrow{#1}}
\newcommand{\cefaout}{\mathcal{O}}
\newcommand{\Lang}{\mathcal{L}}
\newcommand{\Tran}{\mathcal{T}}
\newcommand{\arrseparator}{\sharp}
\newcommand{\todo}[1]{{\color{orange}\textbf{TODO:} #1 \textbf{:ODOT}\color{black}}}
\newcommand{\commentword}[1]{}

\newcommand{\mysplit}{\mathsf{split}}
\newcommand{\mysplitstr}{\mathsf{splitstr}}

\newcommand{\myjoin}{\mathsf{join}}
\newcommand{\mylen}{\mathsf{len}}
\newcommand{\mystrlen}{\mathsf{strlen}}
\newcommand{\myseqlen}{\mathsf{seqlen}}
\newcommand{\myread}{\mathsf{read}}
\newcommand{\mywrite}{\mathsf{write}}
\newcommand{\mymatch}{\mathsf{match}}

\newcommand{\mymatchall}{\mathsf{matchAll}}
\newcommand{\mymatchallstr}{\mathsf{matchAllstr}}


\newcommand{\myat}{{\sf nth}}
\newcommand{\myfilter}{{\sf filter}}
\newcommand{\myset}[1]{\{#1\}}
\newcommand{\hide}[1]{ }
\newcommand{\arraylogic}{{\sf SeqStr}}
\newcommand{\slarraylogic}{{\sf SeqStr}_{\sf SL}}
\newcommand{\exstrlogic}{{\sf XStr}}
\newcommand{\slexstrlogic}{{\sf XStr}_{\sf SL}}

\newcommand*{\inttype}{{\sf Int}}
\newcommand*{\strtype}{{\sf Str}}
\newcommand*{\seqtype}{{\sf Seq}}

\newcommand{\svarx}{\mathfrak{x}}
\newcommand{\svary}{\mathfrak{y}}
\newcommand{\svarz}{\mathfrak{z}}

\newcommand{\svaru}{\mathfrak{u}}
\newcommand{\svarv}{\mathfrak{v}}
\newcommand{\svarm}{\mathfrak{m}}
\newcommand{\svarn}{\mathfrak{n}}
\newcommand{\svars}{\mathfrak{s}}
\newcommand{\svart}{\mathfrak{t}}

\newcommand{\ivarx}{x}
\newcommand{\ivary}{y}
\newcommand{\ivarz}{z}

\newcommand{\seqvara}{\alpha}
\newcommand{\seqvarb}{\beta}
\newcommand{\seqvarg}{\gamma}
\newcommand{\concat}{\cdot}
\newcommand{\eqdef}{\stackrel{\mbox{\begin{tiny}def\end{tiny}}}{=}}


\newcommand{\iop}{\bowtie}
\newcommand{\separator}{\dag}
\newcommand{\enc}{{\sf enc}}
\newcommand{\subseq}{{\sf subseq}}
\newcommand{\elem}{{\sf elem}}

\newcommand*{\product}{\times}

\newcommand{\natnum}{\mathbb{N}}
\newcommand{\intnum}{\mathbb{Z}}

\newcommand{\pref}{{\sf pref}}

\newcommand{\nfa}{{\sf NFA}}
\newcommand{\nft}{{\sf NFT}}


% \newcommand{\oursolver}{{\sf OSTRICH$_{\sf SEQ}$}}


\newcommand{\ostrichseq} {{\sf OSTRICH^{SEQ}}}
\newcommand{\princessarr} {{\sf Princess^{ARR}}}
\newcommand{\princess}{{\sf Princess}}
\newcommand{\seco}{{\sf SeCo}}
\newcommand{\ostrich} {{\sf OSTRICH}}
\newcommand{\seqbase} {{\sf SEQBASE}}
\newcommand{\seqext} {{\sf SEQEXT}}
\newcommand{\withmatch} {{\seqext\mbox{\sf -M}}}
\newcommand{\withoutmatch} {{\seqext\mbox{\sf -N}}}


%\theoremstyle{remark}
%\newtheorem{definition}{Definition}[section]

\newcommand{\zhilin}[1]{\color{blue} {ZL: #1 :LZ} \color{black} }
\newcommand{\denghang}[1]{\color{red} {DH: #1 :HD} \color{black} }
\newcommand{\tl}[1]{\color{magenta} {TL: #1 :LT} \color{black} }
\newcommand{\fu}[1]{\color{purple} {FU: #1 :UF} \color{black} }
\DeclareRobustCommand{\polished}[1]{{\color{red}{#1}}}

%\newcommand{\denghang}[1]{{\color{teal}\textbf{DH:} #1 \textbf{:HD}\color{black}}}
%\newcommand{\comment}[1]{{\color{blue}\textbf{MARK:} #1 \textbf{:KRAM}\color{black}}}


\pagestyle{plain}

%%%%%%%%%%%%%%%%%%%%%%%%%%%%%%%%%%%%%%%%%%%%%%%%%%%%%%%%%%%%%%%%%%%%%%%%%%%%%%%%%%%%%%%%%%

\title{Decision Procedure for \\ A Theory of String Sequences\thanks{This work was partially supported by TBD.}}

% \title{Contribution Title}
% %
% %\titlerunning{Abbreviated paper title}
% % If the paper title is too long for the running head, you can set
% % an abbreviated paper title here
% %
% \author{First Author\inst{1}\orcidID{0000-1111-2222-3333} \and
% Second Author\inst{2,3}\orcidID{1111-2222-3333-4444} \and
% Third Author\inst{3}\orcidID{2222--3333-4444-5555}}
% %
% \authorrunning{F. Author et al.}
% % First names are abbreviated in the running head.
% % If there are more than two authors, 'et al.' is used.
% %
% \institute{Princeton University, Princeton NJ 08544, USA \and
% Springer Heidelberg, Tiergartenstr. 17, 69121 Heidelberg, Germany
% \email{lncs@springer.com}\\
% \url{http://www.springer.com/gp/computer-science/lncs} \and
% ABC Institute, Rupert-Karls-University Heidelberg, Heidelberg, Germany\\
% \email{\{abc,lncs\}@uni-heidelberg.de}}

\author{Denghang Hu\inst{1,2} \and 
  Taolue Chen\inst{3} \and
  Philipp R\"{u}mmer\inst{4}\and\\ 
  Fu Song\inst{1,2,5} \and
  Zhilin Wu\inst{1,2} 
  }

\institute{
Key Laboratory of System Software and State Key
  Laboratory of Computer Science, Institute of Software, Chinese Academy
  of Sciences 
  \and
University of Chinese Academy of Sciences, China
\and 
Birkbeck, University of London, United Kingdom
\and University of Regensburg, Germany
\and Nanjing Institute of Software Technology, China
}


%\author{Denghang Hu\inst{1,2}, Taolue Chen\inst{3}, Philipp R\"{u}mmer\inst{4}, Zhilin Wu\inst{1,2}}
%\institute{Key Laboratory of System Software and State Key Laboratory of Computer Science, Institute of Software, Chinese Academy of Sciences, China \and
%University of Chinese Academy of Sciences, China \and 
%Birkbeck, University of London, United Kingdom \and 
%University of Regensburg, Germany
%}
%%%%%%%%%%%%%%%%%%%%%%%%%%%%%%%%%%%%%%%%%%%%%%%%%%%%%%%%%%%%%%%%%%%%%%%%%%%%%%%%%%%%%%%%%%%

\begin{document}
\maketitle

%%%%%%%%%%%%%%%%%%%%%%%%%%%%%%%%%%%%%%%%%%%%%%%%%%%%%%%%%%%%%%%%%%%%%%%%%%%%%%%%%%%%%%%%%%%

\begin{abstract}
The theory of sequences, supported by many SMT solvers, can model program data types including bounded arrays and lists. Sequences are parameterized by the element data type and provide operations such as accessing elements, concatenation, forming sub-sequences and updating elements. Strings and sequences are intimately related; many operations, e.g., matching a string according to a regular expression, 
splitting strings, or joining strings in a sequence, are frequently used in string-manipulating programs. Nevertheless, these operations are typically not directly supported by existing SMT solvers, which instead only consider the generic theory of sequences.  
In this paper, we propose a theory of string sequences and study its
satisfiability. We show that, while it is undecidable in general,  the
decidability can be recovered by restricting to the straight-line
fragment. This is shown by encoding each string sequence as a string,
and each string sequence operation as a corresponding string operation. We provide pre-image computation for the resulting string operations with respect to automata, effectively casting it into the generic OSTRICH string constraint solving framework. 
We implement the new decision procedure as a tool $\ostrichseq$, and carry out experiments on benchmark constraints generated from \polished{real-world JavaScript programs, hand-crafted templates and unit tests}. The experiments confirm the efficacy of our approach. 
\end{abstract}

%%%%%%%%%%%%%%%%%%%%%%%%%%%%%%%%%%%%%%%%%%%%%%%%%%%%%%%%%%%%

%\vspace{-6mm}
\section{Introduction} \label{sect:intro}

%!TEX root = main.tex

Many real-world applications, such as web applications and programs processing data of type string, involve a multitude of complex operations that convert between strings and string sequences. Typical examples include $\myjoin$ and $\mysplit$, which are present in 
almost all built-in libraries of modern programming languages.
%
Reasoning about string sequences, also known as extendable string arrays~\cite{JezLMR23}, is therefore a crucial aspect of program analysis. 

Satisfiability Modulo Theory (SMT~\cite{smt_18}) solving provides an automatic way to verify software systems. 
%using sequences. 
At the moment, there is no standardized sequence theory in SMT-LIB~\cite{smt_standard}. Instead, in verification tools, often the theory of arrays is used, which is limited and does not provide even basic operations like sequence length. Fortunately, \cite{seq_2012} proposed a basic format for a sequence theory, which has been implemented and extended by SMT solvers such as Z3~\cite{Z3} and cvc5~\cite{BarbosaBBKLMMMN22} with practical operations. In particular, \cite{cvc5_seq} presents a calculus in cvc5 for reasoning about string sequences, which extends the proposal~\cite{seq_2012} with \verb|update| (aka \verb|write|) operations. 
The implementation of sequences in Z3 has been documented in the Z3 guide~\cite{z3_seq_doc} and includes operations such as \verb|map| and \verb|foldleft| (but no \verb|update|).

Although existing SMT solvers are effective at reasoning about sequences in general, to the best of our knowledge, none of them directly support functions that convert between strings and string sequences %.  no SMT solvers can handle conversions between string and string sequence, for instance through 
such as $\myjoin$ and $\mysplit$. While it is possible to implement these functions using existing ones for some SMT solvers (e.g., Z3 can use \verb|map| and \verb|foldleft| to implement the $\myjoin$ function), usually the solvers do not provide decision procedures for those theories and fail to solve many instances. 
This limitation poses a significant challenge for program analysis that involves such operations.
%
% illustrate the need for a new solver
The purpose of our work is to fill this gap. 
%between the demand from program analysis of real-world string-manipulating programs and the capabilities of the state-of-the-art SMT solvers.

%%%%%%%%%%%%%%%%%%%%%%%%%%%%%%%%%%%%%%%%%%%%%%%%%%%%%%%

\smallskip
\noindent\emph{Contributions.} We propose a logic (aka constraint language) $\arraylogic$ of string sequences, which extends the existing sequence theory in SMT solvers with dedicated operations of string sequences and strings. In addition to the standard string operations (replace/replaceAll, reverse, finite transducers, regular constraints, 
%as well as 
length, indexOf, substring, etc.), our logic emphasizes the interaction between strings and string sequences through operations such as $\mysplit$, $\myjoin$, $\myfilter$, and $\mymatchall$. For instance, $\mysplit$ splits a string into a sequence of substrings according to a regular expression; $\myjoin$ concatenates all the strings in a sequence to obtain a single string. More details can be found in Section~\ref{sec:logic}. Typically, when string sequences are considered, it is natural to consider the integer data type as well; for instance, one usually needs to refer to the length of a sequence or the index of a specific sequence element. As a result, $\arraylogic$ features three data types: integers, strings, and string sequences. In terms of operations, the logic subsumes, to our best knowledge, most string constraint languages that have been proposed in the literature. 

%
%This motivates the following problem: provide a decision procedure that supports both string and integer data type, with completeness guarantee and meanwhile
%admitting efficient implementation.

Our main focus is the decision procedure for the satisfiability problem of  $\arraylogic$. 
Because of its expressiveness, it is perhaps not surprising that $\arraylogic$ is undecidable in general. To reinstate decidability, we consider the straight-line fragment, which imposes a syntactic restriction on the constraints written in the logic. %In a nutshell, this
Straight-line formulas naturally arise when verifying programs using
bounded model checking or  symbolic execution, 
% (e.g. see [27, 43, 65]), 
which unroll loops in the programs up to a given
depth and convert programs to static single assignment form (i.e. each variable is defined at most once). For example, a majority of the constraints in the standard Kaluza
benchmarks~\cite{Berkeley-JavaScript} %with 50,000+ test cases generated by symbolic execution of
%JavaScript applications 
satisfy this condition. (Cf.\ Section~\ref{sect:example} for a concrete example.)  


%Nonetheless, as already noted in
%[Lin and Barceló 2016], the straight-line restriction is reasonable since it is typically satisfied by
%constraints that are generated by symbolic execution, e.g., all constraints in the standard Kaluza
%benchmarks [Saxena et al. 2010] with 50,000+ test cases generated by symbolic execution on
%JavaScript applications were noted in [Ganesh et al. 2012] to satisfy this condition. Intuitively,
%as elegantly described in [Bjørner et al. 2009], constraints from symbolic execution on stringmanipulating programs can be viewed as the problem of path feasibility over loopless stringmanipulating programs S with variable assignments and assertions

%Fortunately, recent years have seen the possibility of recovering some decidability of string constraint languages with multiple string operations, while retaining applicability for constraints that
%arise in practical symbolic execution applications. This is done by imposing syntactic restrictions including acyclicity [Abdulla et al. 2014; Barceló et al. 2013], solved form [Ganesh et al. 2012],
%and straight-line [Chen et al. 2018a; Holík et al. 2018; Lin and Barceló 2016]. These restrictions
%are known to be satisfied by many existing string constraint benchmarks, e.g., Kaluza [Saxena
%et al. 2010], Stranger [Yu et al. 2010], SLOG [Holík et al. 2018; Wang et al. 2016], and mutation XSS
%benchmarks of [Lin and Barceló 2016].

%The straight-line logic of [Lin
%and Barceló 2016] unified the earlier logics by allowing concatenation, regular constraints, rational
%transductions, and length and letter-counting functions. It was pointed out by [Chen et al. 2018a]
%that this logic cannot express the replaceAll function with the replacement string provided as a
%variable, which was never studied in the context of verification and program analysis. Chen et al.
%proceeded by showing that a new straight-line logic with the more general replaceAll function and
%concatenation is decidable, but becomes undecidable when the length function is permitted.

The general strategy of our decision procedure is to encode each string sequence as a string in such a way that all operations for string sequences considered in $\arraylogic$ can be transformed into string operations, possibly involving integers. As such, the decidability of (straight-line) $\arraylogic$ is reduced to string constraints with the integer data type, which was shown to be decidable \cite{atva2020}. Needless to say, this strategy requires overcoming certain technical challenges, as the string constraints resulting from the reduction typically encompass complex (and sometimes non-standard) string operations, which go well beyond the capacity of state-of-the-art string constraint solvers. % to solve, since they usually involve the string operations {\tt str.replace\_re\_all}, {\tt str.substr}, {\tt str.len}, as well as some non-standard  
Recall that for the decidability of string constraints with integers, cost-enriched finite automata (CEFA), a variant of cost-register automata,  %introduced by Alur et al. [4], which are called  for convenience. 
were utilized \cite{atva2020}. %to maintain a separation of string variables and integer variables. 
%Intuitively, each CEFA records the connection between a
%string variable and its associated integer variables. With CEFAs, the concept of recognisable relations is then naturally extended to accommodate integers. %The integer constraints, however, are detached from CEFAs rather than being part of CEFAs. 
%This allows to preserve the independence of string variables in the recognisable relation. 
The crux of our technical contributions is thus to compute the backward images of CEFAs under the fairly complex string operations (cf. Section~\ref{subsect:preimage}), %inherited from %, where each cost register (integer variable) might be split into several ones. %, thus 
%The advantage of this approach lies in extending but still in the same 
so the powerful string constraint solving framework  OSTRICH %for string constraints
%without the integer data type 
\cite{CCH+18,CHL+19} can be harnessed. %, which is able to treat a wide range of string operations in a generic way. 
%To the best of our knowledge, the
%%class of string constraints considered in this paper is currently one of the most expressive string theories involving the integer data type known to enjoy a decision procedure.

%-------------------------------------------------------------------------------------
We implement the decision procedure as a new solver $\ostrichseq$ on top of OSTRICH~\cite{CHL+19} and $\princess$~\cite{princess08}. To evaluate the effectiveness of $\ostrichseq$, we curate two benchmark suites of constraints that 
are randomly generated from templates (only using
the operations directly supported by some existing SMT solvers) and extracted from real-world JavaScript programs and unit tests (involving string sequence operations that are not directly supported by existing SMT solvers), respectively. 
%generated from real-world and hand-crafted JavaScript programs, 
%We compare $\ostrichseq$ with the SOTA string solvers, cvc5, Z3, Z3-noodler~\cite{ChenCHHLS24}, $\ostrich$ and $\princessarr$~\cite{princess08}.
%
On both benchmark suites, the experimental results show that $\ostrichseq$ can solve considerably more constraints
%generated from both real-world and hand-crafted JavaScript programs 
than SOTA string solvers including cvc5, Z3, Z3-noodler~\cite{ChenCHHLS24}, $\ostrich$ and $\princessarr$~\cite{ruemmer2024princess}, while being largely as efficient as most of them. 

%The implementation consists of approximately 5,000 lines of Scala code. The core functionality—preimage computation for sequence operations such as $\mysplit$, $\myjoin$, $\myfilter$, and $\mymatchall$—is implemented in about 2,000 lines of code and is integrated directly into the OSTRICH framework. To support the representation of sequences, we built a lightweight library for sequence automata that compiles the sequence variable into a set of automata, which are then used to solve the constraints. The library is implemented in about 1,000 lines of code. The remaining code is used for the integration with OSTRICH and for the implementation of the decision procedures for the theory of sequences. 

%------------------------------------------------------------------------
\smallskip
\noindent\emph{Organization.} %The rest of the paper is structured as follows: 
Section~\ref{sect:example} presents a motivating example. Section~\ref{sec:prelim} gives the preliminaries. 
Section~\ref{sec:logic} defines the logic $\arraylogic$ of string sequences. Section~\ref{sect:procedure} presents the decision procedure. Section~\ref{sect:exp} presents the benchmarks and
experiments. 
Section~\ref{sect:related} discusses the related work. The paper is concluded in Section~\ref{sect:conc}.




%\vspace{-2mm}
\section{Preliminaries}\label{sec:prelim}

%!TEX root = main.tex


%Throughout the paper, 
Let $\natnum$ denote the set of natural numbers. For $1\leq n\in \natnum$, let $[n]:=\{1, \ldots, n\}$, and for $m<n \in \natnum$, let $[m,n] := \{j \mid m \le j \le n\}$. We also use standard quantifier-free/existential \emph{linear integer arithmetic} (LIA) formulas, which are typically ranged over by $\phi, \varphi$, etc. For an integer term $t$, we use $t[t_1/t_2]$ to denote the term obtained by replacing $t_2$ with $t_1$ where $t_1$ and $t_2$ are integer terms.

\smallskip
\polished{\noindent \textbf{Strings, languages, and transductions.}} We fix an alphabet $\Sigma$, i.e., a finite set of letters. 
%Moreover, we use $\nullchar$ to denote a special null symbol and assume that  $\nullchar \not \in \Sigma$. 
A \emph{string} over $\Sigma$ is a finite sequence of letters from $\Sigma$. We use $\Sigma^*$ to denote the set of strings over $\Sigma$ and $\varepsilon$ to denote the empty string. %Moreover, for convenience, we use $\Sigma^\varepsilon$ to denote $\Sigma \cup \{\varepsilon\}$. 
%A string $w'$ is called a \emph{prefix} of $w$ if $w = w'w''$ for some string $w''$. We use $\pref(w)$ to denote the set of prefixes of $w$. For a prefix $w_1$ of $w$, let $w = w_1 w_2$, then we use $w_1^{-1}w$ to denote $w_2$.
%
A \emph{language} over $\Sigma$ is a subset of $\Sigma^*$. We will use $L_1, L_2, \dots$ to denote languages. For two languages $L_1, L_2\subseteq \Sigma^*$, $L_1 \cup L_2$ denotes the union of $L_1$ and $L_2$, and $L_1 \cdot L_2$  denotes the concatenation of $L_1$ and $L_2$, that is,  $\{u_1 \cdot u_2 \mid u_1 \in L_1, u_2 \in L_2\}$. For a language $L\subseteq\Sigma^*$, we define the complement of $L$ as $\Bar{L} = \{w\in\Sigma^*\mid w\not\in L\}$, moreover, we define $L^n$ for $n \in \natnum$, the \emph{iteration} of $L$ for $n$ times, inductively as: $L^0=\{\varepsilon\}$ and $L^{n} =L \cdot L^{n-1}$ for $n > 0$. We also use $L^*$ to denote the iteration of $L$ for arbitrarily many times, that is, $L^* = \bigcup_{n \in \natnum} L^n$, and let $L^+ = \bigcup_{n>0} L^n$. \polished{A \emph{transduction} over $\Sigma$ is a binary relation over $\Sigma^*$, namely, a subset of $\Sigma^* \times \Sigma^*$.}
%\begin{definition}[

Regular expressions over $\Sigma$ are defined in a standard way, i.e., 
$$e \eqdef \emptyset \mid \varepsilon \mid a \mid e + e \mid e \concat e \mid e^*  \text{where}\  a \in \Sigma.$$
%

The language $\Lang(e)\subseteq\Sigma^*$ of a regular expression $e$ is defined %that is, the set of strings that match $e$, 
	inductively as: $\Lang(\emptyset) =\emptyset$,
	$\Lang(\varepsilon) =\{\varepsilon\}$,
	$\Lang(a)= \{a\}$,
	$\Lang(e_1 + e_2) = \Lang(e_1) \cup \Lang(e_2)$,
	$\Lang(e_1 \concat e_2) = \Lang(e_1) \cdot \Lang(e_2)$,
	$\Lang(e_1^*)=(\Lang(e_1))^*$. \polished{A \emph{regular language} is a language that can be defined by a regular expression.}
%\end{definition}


\hide{As $+$ is associative and commutative, we write $(e_1 + e_2) + e_3$ as $e_1 + e_2 + e_3$ for brevity. We use the abbreviation $e^+ \equiv e \concat e^*$. Moreover, for $\Gamma = \{a_1, \cdots, a_n\}\subseteq \Sigma$, we may overload $\Gamma$ (resp. $\Gamma^\ast$) as the regular expression $a_1 + \cdots + a_n$ (resp. $(a_1 + \cdots + a_n)^\ast$). 
In addition, $|e|$ denotes the number of symbols occurring in $e$.}

\smallskip
\noindent \textbf{Automata.} A \emph{(nondeterministic) finite automaton} (NFA) $\NFA$ is a 5-tuple $(Q, \Sigma, \delta,$ $ I, F)$, where $Q$ is a finite set of states, $\Sigma$ is a finite alphabet, $\delta \subseteq Q \times \Sigma \times Q$ is the transition relation, $I,F \subseteq Q$ are the sets of initial and final states respectively. For readability, we write a transition $(q, a, q') \in \delta$ as $q \xrightarrow[\delta]{a} q'$ (or simply $q \xrightarrow{a} q'$). %Moreover, when $\delta$ is clear from context, we omit $\delta$ in $q \xrightarrow[\delta]{a} q'$ and write $q \xrightarrow{a} q'$. 
%
A \emph{run} of $\NFA$ on a string $w = a_1 \cdots a_n$ is a sequence of transitions $q_0 \xrightarrow{a_1} q_1 \cdots q_{n-1} \xrightarrow{a_n} q_n$ with $q_0 \in I$. The run is \emph{accepting} if $q_n \in F$.
A string $w$ is accepted by an NFA $\NFA$ if there is an accepting run of $\NFA$ on $w$. In particular, the empty string $\varepsilon$ is accepted by $\NFA$ if $I \cap F \neq \emptyset$. The language of $\NFA$, denoted by $\Lang(\NFA)$, is the set of strings accepted by $\NFA$. In addition, an NFA can be built to recognize the language of each given regular expression, and vice versa.
%the language of an NFA corresponds 
%to the language described by a specific regular expression.  

A \emph{(nondeterministic) finite transducer (\nft)} $\transducer$ is an extension of {\nfa} with outputs. Formally, an {\nft} $\transducer$ is a 5-tuple $(Q, \Sigma, \delta, I, F)$, where $Q, \Sigma, I, F$ are the same as in {\nfa} and the transition relation $\delta$ is a finite subset of $Q \times \Sigma \times Q \times \Sigma^*$. For readability, we write a transition $(q, a, q', u) \in \delta$ as $q \xrightarrow[\delta]{a, u} q'$ or $q \xrightarrow{a, u} q'$. A run of $\transducer$ over a string $w=a_1 \cdots a_n$ is a sequence of transitions $q_0 \xrightarrow{a_1, u_1} q_1\xrightarrow{a_2, u_2} q_2\cdots q_{n-1} \xrightarrow{a_n, u_n} q_n$
where $q_0 \in I$. Similar to {\nfa}, the run is  \emph{accepting} if $q_n \in F$. The string $u_1 \cdots u_n$ is called the \emph{output} of the run. The transduction $\Tran(\transducer)\subseteq \Sigma^*\times \Sigma^*$ defined by $\transducer$ is the set of string pairs $(w, u)$ such that there is an accepting run of $\transducer$ on $w$, with the output $u$.
\begin{definition}[Cost-enriched finite automata]
A cost-enriched finite automaton (CEFA for short) $\aut$ is a 6-tuple $(R, Q, \Sigma, \delta, $ $I, F)$ where
	\begin{itemize}[topsep=0pt,partopsep=0pt]
            \item $R = \{r_1, \cdots, r_k\}$ is a finite set of registers,
		\item $Q, I, F$ are the same as in {\nfa}, i.e., the set of states, the set of initial states, and the set of final states, respectively, 
		\item $\delta \subseteq Q \times \Sigma \times Q \times \Int^{k}$ is a transition relation, where $\Int^{k}$ represents the values to update the registers in $R$. 
	\end{itemize}
We write $R_\aut$ for the set of registers of $\aut$
and represent it as a vector $(r_1, \cdots, r_k)$. Accordingly, updates $r_i:= r_i + v_i$ for all $i\in [k]$  are simply identified as the vector $\vec{v}=(v_1, \cdots, v_k)$, i.e., $r_i$ is  incremented by $v_i$ for each $i \in [k]$. Typically, we write a transition $(q, a, q', \vec{v}) \in \delta$ as $q \xrightarrow[\vec{v}]{a} q'$.
\end{definition}    
%\hide{In particular, $(q,a,q')$ stands for the transition $(q,a,q',())$ without any updates.}
Intuitively, CEFAs add write-only cost registers to finite automata, where ``write-only'' means that the cost registers can only be written/updated but cannot be read, i.e., they cannot be used in the guards of the transitions.  


%The semantics of CEFA is defined as follows. 
%Let $\aut = (R, Q, \Sigma, \delta, I, F)$ be a CEFA. 
%A \emph{configuration} of $\aut$ is a pair $(q, \vec{v})$ where $q \in Q$ and $\vec{v}$ is a vector denoting the values of registers.  
%An \emph{initial configuration} of $\aut$ is $(q_0, \vec{0})$ with $q_0 \in I$, where the value of each register is zero. 
A \emph{run} of the CEFA $\aut$ on a string $w = a_1 \cdots a_n$ is a sequence of transitions $q_0 \xrightarrow[\myvec{v_1}]{a_1} q_1 \cdots q_{n-1}\xrightarrow[\myvec{v_n}]{a_n} q_n$ such that $q_0 \in I$ and $q_{i-1} \xrightarrow[\myvec{v_i}]{a_i} q_i$ for each $i \in [n]$. 
%
%The run $q_0 \xrightarrow[\myvec{v_1}]{a_1} q_1 \cdots q_{n-1}\xrightarrow[\myvec{v_n}]{a_n} q_n$ 
It is \emph{accepting} if $q_n \in F$, and 
%
the vector $\myvec{c}= \sum_{i \in [n]} \myvec{v_i}$ is defined as the \emph{cost} of the run. %$q_0 \xrightarrow[\myvec{v_1}]{a_1} q_1 \cdots q_{n-1}\xrightarrow[\myvec{v_n}]{a_n} q_n$. 
(Note that all registers are initialized to zero.) 
%
%and updated to $\sum \limits_{j \in [n]} \myvec{v_j}$ after all the transitions in the run are executed. 
We write $\myvec{c} \in \aut(w)$ if there is an accepting run of $\aut$ on $w$ whose cost is $\myvec{c}$.  The language of a CEFA $\aut$, denoted by $\Lang(\aut)$, is defined as the set of pairs $\{(w, \myvec{c})\in\Sigma^*\times \Int^{|R|} \mid  \myvec{c} \in \aut(w)\}$.
%
In particular, if $I \cap F \neq \emptyset$, then $(\varepsilon, \myvec{0}) \in \Lang(\aut)$. 
We denote by $\Lang_1(\aut)$ the language $\{w\in\Sigma^*\mid \exists \myvec{c}\in \Int^{|R|}.\ (w, \myvec{c})\in \Lang(\aut)\}$.

%\hide{Moreover, the \emph{output} of a CEFA $\aut$, denoted by $\cefaout(\aut)$, is defined as $\{\myvec{c} \mid  \exists w.\ \myvec{c} \in \aut(w)\}$.
%}



%%%%%%%%%%%%%%%%%%%%%%%%%%%%%%%%%%%%%%%%%%%%%%%%%%%%%%%%%%%%%%%%%%%%%%%%%%%%%%%%%
\hide{The following example illustrates the backward propagation. 

% \vspace*{-2mm}
\begin{example}
	Consider the equality $\svarv = f(\svaru_1, it_1, \svaru_2, it_2)$ (where $\svaru_1,\svaru_2$ are string variables, $it_1, it_2$ are integer terms, and $f$ is a function) and the CEFA $\aut = (R, Q, \Sigma, \delta, I, F)$, where $R= \{r_1\}$. %(that is, $R$ contains only one register). 
	%Moreover, let $t = r^{(1)}_1 + r^{(2)}_1$ which intuitively describes how $r_1$ can be calculated from $r^{(1)}_1$ and $r^{(2)}_1$. 
	Then an $R$-cost-enriched pre-image of $\aut$ under $f$, denoted by $f^{-1}_{R}(\aut)$, is a pair $(\cerl, t)$ where $\cerl \subseteq (\Sigma)^\ast \times (\intnum \times \intnum) \times (\Sigma)^\ast \times (\intnum \times \intnum)$, $t$ is an LIA term over $\{r^{(1)}_1, r^{(2)}_1\}$, and $\cerl$ satisfies that for every $(v, n') \in \Lang(\aut)$, there is $(u_1, n_1, n'_1, u_2, n_2, n'_2) \in \cerl$ such that $f(u_1, n_1, u_2, n_2) = v$, and  $n'= t[n'_1/r^{(1)}_1, n'_2/r^{(2)}_1]$. 
	%
	We say that a pre-image $f^{-1}_{R}(\aut) = (\cerl, t)$ is CEFA-definable if $\cerl$ can be expressed as a finite collection of CEFA tuples $(\autb_{j, 1}, \autb_{j, 2})$ with $j \in [k]$ such that $R_{\autb_{j,1}}=(r'_{1,1}, r^{(1)}_1)$, $R_{\autb_{j, 2}} = (r'_{2,1}, r^{(2)}_1)$, and $\cerl = \bigcup_{j \in [k]} \Lang(\autb_{j, 1}) \times \Lang(\autb_{j, 2})$.  
	Intuitively, $r'_{1,1}$ and $r'_{2,1}$ are the fresh registers corresponding to the two integer parameters, and $r^{(1)}_1, r^{(2)}_1$ are the fresh registers introduced in $\autb_{j, 1}$ and $\autb_{j, 2}$ respectively for computing $r_1$ according to $t$. Then the backward propagation of the constraint $\svarv \in \aut$ 
	%with respect to the equality $\svarv = f(\svaru_1, it_1, \svaru_2, it_2)$ 
	under $f$ nondeterministically chooses a tuple $(\autb_{j, 1}, \autb_{j, 2})$, removes $\svarv = f(\svaru_1, it_1, \svaru_2, it_2)$ from the string constraint, and adds $\svaru_1 \in \autb_{j, 1} \wedge \svaru_2 \in \autb_{j, 2} \wedge r_1 = t \wedge it_1 = r'_{1,1} \wedge it_2 = r'_{2,1}$ to the string constraint.
	%removes $\svarv = f(\svaru_1, it_1, \svaru_2, it_2)$ from the string constraint, and adds $r_1 = t \wedge it_1 = r'_{1,1} \wedge it_2 = r'_{2,1}$ to the string constraint. \qed
\end{example}
% \vspace*{-2mm}
}

\hide{
\begin{definition}[Finite-state automata] \label{def:nfa}
	A \emph{(nondeterministic) finite-state automaton
	(\nfa)} over the alphabet $\Sigma$ is a tuple $\aut =
	(Q, \Sigma, \delta, I, F)$, where 
	$Q$ is a finite set of 
	states, $I \subseteq Q$ is
	a set of initial states, $F \subseteq Q$ is a set of final states, and 
	$\delta \subseteq Q \times \Sigma \times Q$ is the
	transition relation. 
\end{definition}

Given an input string $w\in \Sigma^*$, a \emph{run} of $\aut$ on $w$
%(with $a_0 = \EndLeft$ and $a_{n+1} = \EndRight$)
is a sequence $q_0 a_1 q_1 \ldots a_n q_n$ such that $q_0 \in I$, $w = a_1 \cdots a_n$ and $(q_{j-1}, a_{j}, q_{j}) \in
\delta$ for every $j \in [n]$.
The run is \emph{accepting} if $q_n \in F$.
A string $w$ is \emph{accepted} by $\aut$ if there is an accepting run of
$\aut$ on $w$. In particular, the empty string $\varepsilon$ is accepted by $\aut$ if $I \cap F \neq \emptyset$. 
The set of strings accepted by $\aut$, i.e., the language \emph{recognized} by $\aut$, is denoted by $\Lang(\aut)$.
%Since we deal with computational complexity in the sequel, we define
The \emph{size} $|\aut|$ of $\aut$ is defined as the cardinality of the set $Q$ of states, which will be 
used when the computational complexity is concerned.

For convenience, $(q,a,q')\in\delta$ may be written as  $q \xrightarrow[\delta]{a} q'$ or $q \xrightarrow{a} q'$,
and for each letter $a \in \Sigma$, let $\delta^{(a)}$  denote  the  relation $\{(q, q') \mid (q, a, q') \in \delta\}$.
}

\polished{
\smallskip
\noindent \textbf{Pre-images under string functions.} Consider a language $L \subseteq \Sigma^* \times \Int^{k_0}$ defined by a CEFA $\aut =(R, Q, \Sigma, \delta, I, F)$ with the registers $R= (r_1, \cdots, r_{k_0})$ and a function $f: (\Sigma^* \times \Int^{k_1}) \times \cdots \times (\Sigma^* \times \Int^{k_l}) \rightarrow \Sigma^*$.  
For each $i \in [k_0]$, let $r^{i}_1, \cdots, r^{i}_l$ be the freshly introduced registers,
and 
$\vec{t} = (t_1, \cdots ,t_{k_0})$ be a vector of LIA formulas such that $t_i$ is a linear combination of $r^{i}_1, \cdots, r^{i}_l$ for each $i \in [k_0]$. 

The pre-image of $L$ under $f$ with respect to $\vec{t}$, denoted by $f^{-1}_{\vec{t}}(L)$, is 
%$\cerl$ such that
%\begin{itemize}
%\item $\vec{t} = (t_1, \cdots ,t_{k_0})$ is a vector of LIA formulas where $t_i$ is a linear combination of fresh registers $r^{i}_1, \cdots, r^{i}_l$ for each $i \in [k_0]$;
%
%[intuitively, each cost register $r_i$ is split into $l$ cost registers $r^{(1)}_i, \cdots,r^{(l)}_i$, one for each string input parameter, and $t_i$ tells how to compute $r_i$ from $r^{(1)}_i, \cdots,r^{(l)}_i$]
%\item %$\cerl\subseteq (\Sigma^* \times \Int^{k_1 + k_0}) \times \cdots \times (\Sigma^* \times \Int^{k_l + k_0})$ such that 
a relation 
\begin{center}
    $\cerl\subseteq  (\Sigma^* \times \Int^{k_1 + k_0}) \times \cdots \times (\Sigma^* \times \Int^{k_l + k_0})$
\end{center}
that comprises tuples of the form  
$((w_1, (\myvec{c_1}, \myvec{d_1})), \cdots, (w_l, (\myvec{c_l}, \myvec{d_l})))$ 
such that
\begin{itemize}
\item 
for every $j\in [l]$,
$\myvec{c_j} \in \Int^{k_j}$ and $\myvec{d_j} = (d^{1}_{j}, \cdots, d^{k_0}_{j}) \in \Int^{k_0}$, 
%
\item let 
$w_0 = f((w_1, (\myvec{c_1}, \myvec{d_1})), \cdots, (w_l, (\myvec{c_l}, \myvec{d_l})))$ and  
\begin{center}
    $\myvec{d'} = \left(t_1\left[d^{1}_{1}/r^{1}_1, \cdots, d^{1}_{l}/r^{1}_l\right], \cdots, t_{k_0}\left[d^{k_0}_{1}/r^{k_0}_{1}, \cdots, d^{k_0}_{l}/r^{k_0}_{l}\right] \right),$
\end{center}
it holds that 
$(w_0, \myvec{d'}) \in L$.
\end{itemize}
%
%ensure that $L$ can be derived from $\cerl$ and $\vec{t}$, and comprising
%
%\[\left(w_0, t_1\left[d^{1}_{1}/r^{1}_1, \cdots, d^{1}_{l}/r^{1}_l\right], \cdots, t_{k_0}\left[d^{k_0}_{1}/r^{k_0}_{1}, \cdots, d^{k_0}_{l}/r^{k_0}_{l}\right]\right), \]
%
%such that  
%\[w_0 = f\left((w_1, \vec{c_1}), \cdots, (w_l, \vec{c_l}\right)) \mbox{ for some } ((w_1, (\vec{c_1}, \vec{d_1})), \cdots, (w_l, (\vec{c_l}, \vec{d_l}))) \in \cerl,\]
%where $\vec{c_j} \in \Int^{k_j}$, $\vec{d_j} = (d^{1}_{j}, \cdots, d^{k_0}_{j}) \in \Int^{k_0}$ for $j\in [l]$.
%
%$\vec{c_1} \in \Int^{k_1}$, $\cdots$, $\vec{c_l} \in \Int^{k_l}$, $\vec{d_1} = (d^{(1)}_{1}, \cdots, d^{(1)}_{k_0}) \in \Int^{k_0}$, $\cdots$, and $\vec{d_l} = (d^{(l)}_{1},\cdots, d^{(l)}_{k_0}) \in \Int^{k_0}$.
%\end{itemize}

The pre-image $f^{-1}_{\vec{t}}(L)$ is \emph{CEFA-definable} if $f^{-1}_{\vec{t}}(L) = \bigcup_{i=1}^n \Lang(\aut_{i,1}) \times \cdots \times \Lang(\aut_{i,l})$ for some $n\ge 1$, where for all $i\in [n]$ and $j\in [l]$, $\aut_{i,j}$ is a CEFA such that $\Lang(\aut_{i,j})\subseteq \Sigma^* \times \Int^{k_j}$.  
%
Finally, a pre-image of $L$ under $f$, denoted by $f^{-1}(L)$, is a pair $(\cerl,\vec{t})$ such that $\cerl = f^{-1}_{\vec{t}}(L)$.
%is the pre-image of $L$ under $$f$ with respect to $\vec{t}$. 
%

The core concept of \emph{pre-image computation} for the function $f$ involves determining the vector $\vec{t}$ of LIA formulas and, for each CEFA $\aut$, computing $f^{-1}_{\vec{t}}(\Lang(\aut))$ represented as a finite set of CEFAs. 
%
Intuitively, we remove the string equalities of the form $y = f(x_1, \myvec{it_1}, \cdots, x_l, \myvec{it_l})$ (where $\myvec{it_1}$, $\cdots$, $\myvec{it_l}$ are integer expressions) one by one by computing the pre-images and adding the resulting CEFA membership constraints for $x_1, \cdots, x_l$. In the end,  the original string constraint is transformed into a conjunction of CEFA membership constraints and LIA formulas, whose satisfiability checking %can be solved in 
is known to be PSPACE-complete~\cite{atva2020}.
}
 

%An {\nft} $\transducer$ is \emph{deterministic} if $I$ is a singleton, and, for every pair $(q,a) \in Q\times \Sigma$,  there is at most one pair $(q', u) \in Q \times \Sigma^*$ such that $(q, a, q', u) \in \delta$. \tl{do we need this?}
%

%\tl{to put in a different place later} Note that our decision procedure is applicable to general transducers as well, however, the EXPSPACE complexity bound is not, because the distributive property $f^{-1}(L_1\cap L_2)= f^{-1}(L_1)\cap f^{-1}(L_2)$ for regular languages $L_1, L_2$ only holds for functional transducers $f$.  

%We remark that an FT usually defines a relation.

%\smallskip

%In this paper, we consider logics involving two data-types, i.e., the string data-type and the integer data-type. We will use $u, v, \dots$ to denote string constants,  $c, d,\dots$ to denote integer constants, $x, y, \dots$ to denote string variables, and $i, j, \dots$ to denote  integer variables.



%\vspace{-2mm}
\section{Motivating Example} \label{sect:example}

%!TEX root = main.tex

\section{Motivating Example} \label{sect:example}
To motivate and illustrate our approach, we consider a snippet of a Javascript program shown  in Listing~\ref{lst:js_example}. The program is extracted, with slight adaptions, from a real-world example.\footnote{\url{https://github.com/hgoebl/mobile-detect.js/blob/master/mobile-detect.js}} 
%Listing \ref{lst:js_example} is an example of JavaScript code coming from a real-world program\footnote{\url{https://github.com/hgoebl/mobile-detect.js/blob/master/mobile-detect.js}}, which manipulates a sequence of strings. 
%In Listing~\ref{lst:js_example}, 
Here, the function \texttt{prepareVersionNo} extracts from the input string \texttt{version} a version number in the floating point number format. 
Specifically, it splits  \texttt{version} according to the non-numeric characters into a sequence/array {\tt numbers}. If {\tt numbers} contains exactly one element, {\tt numbers[0]} is returned. Otherwise, the function returns the concatenation of {\tt numbers[0]}, the dot, and the string obtained by concatenating the other elements of {\tt numbers}.  For instance, if {\tt version = 12a56b23}, the output %of \texttt{prepareVersionNo} 
is {\tt 12.5623}.

To verify the functional correctness of  \texttt{prepareVersionNo}, we specify the following precondition and postcondition.  
\begin{itemize}
\item The  precondition {\tt version $\in$ \verb|[0-9]+([0-9a-zA-Z._ /-])*|} requires that the input string starts with a decimal number, followed by a string of digits, letters, dot \verb|.|, underline  \verb|_|, blank symbol, slash \verb|/|, or hyphen \verb|-|.
\item As postcondition {\tt result} $\in$ \verb|[0-9]+(\.[0-9]+)?|, the output string starts with a decimal number, possibly followed by a dot \verb|.| and a decimal number. 
\end{itemize}
 
\lstinputlisting[float=t,caption={JavaScript code snippet: The motivating example}, label = {lst:js_example}]
{mot-exmp-js-code.tex}  



\begin{lstlisting}[float=t,language=SMT-LIB, caption={Constraints for path-2 in {\tt prepareVersionNo}}, label={lst:js_example_constraints},belowskip=-2\baselineskip,aboveskip=-0.8\baselineskip]
; precondition 
(assert (str.in_re version preReg)) 
(assert (= numbers (str.splitre version splitReg)))
(assert (< 1 (seq.len numbers)))
(assert (= temp (str.++ (seq.nth numbers 0) ".")))
(assert (= numbers1
           (seq.extract numbers 1 (- (seq.len numbers) 1))))
(assert (= result (str.++ temp (seq.join numbers1 ""))))
; postcondition
(assert (not (str.in_re result postReg))) 
\end{lstlisting} 

We apply symbolic execution to verify \texttt{prepareVersionNo}. % satisfies the specification. 
We enumerate the execution paths of \texttt{prepareVersionNo} and check for each path that, provided the input string satisfies the precondition, after the execution the output string satisfies the postcondition. 
%
%From Listing~\ref{lst:js_example}, we can see that 
There are three execution paths: path-1, taking \verb|numbers.length===1|; path-2, taking \verb|numbers.length>1|; and path-3,  in which both conditions are false.
%Let us exemplify the symbolic execution by considering 
As an example, we consider path-2, where %, that is, the path takes the branch . 
the symbolic execution %of path 2 with respect to the precondition and postcondition is 
reduces to deciding the satisfiability of the SMT formula in Listing~\ref{lst:js_example_constraints}, where {\tt preReg} and {\tt postReg} denote the regular expressions  in the precondition and postcondition respectively, i.e., \verb|[0-9]+([0-9a-zA-Z._ /-])*| and  \verb|[0-9]+(\.[0-9]+)?|, and {\tt splitReg} represents \verb|[a-zA-Z._ /-]|. 

%From Listing~\ref{lst:js_example_constraints}, we can see that 
The SMT formula involves various operations of string sequences and strings, including sequence length {\tt seq.len}, string split {\tt str.splitre}, sequence read {\tt seq.nth}, string concatenation {\tt str.++},
sequence extract {\tt seq.extract}, sequence join {\tt seq.join}, as well as regular constraints  {\tt str.in\_re}. 
While state-of-the-art SMT solvers like Z3 and cvc5 can solve constraints in the generic sequence theory, they do not directly support complex operations that feature mutual transformations between strings and string sequences, such as {\tt str.splitre} and {\tt seq.join}. 
This motivates us to define an SMT theory of string sequences ($\arraylogic$, cf. Section~\ref{sec:logic}) and investigate its decision procedure (Section~\ref{sect:procedure}). %how to solve the constraints in $\arraylogic$. 

\enlargethispage{2ex}%
Due to the general undecidability of %we show that the satisfiability of 
$\arraylogic$, %is undecidable in general. %and it is decidable for the straight-line fragment of $\arraylogic$. 
we focus on %For the decidable 
the straight-line fragment, into which the constraint in Listing~\ref{lst:js_example_constraints} falls.  %The decision procedure of the straight-line fragment is obtained by first 
%
We encode string sequences as strings, and, accordingly, %and 
operations involving string sequences as string operations. %,  then utilize the string solvers to discharge the resulting string constraints. 
%It turns out that the string constraints resulted from the translation are complex and challenging for the SOTA string solvers to solve, since they usually involve the string operations {\tt str.replace\_re\_all}, {\tt str.substr}, {\tt str.len}, as well as some non-standard string operations. 
In the example, the resulting constraint  %from the translation of the constraint in Listing~\ref{lst:js_example_constraints} 
contains a complex combinations of  {\tt str.replace\_re\_all}, {\tt str.substr} and {\tt str.len} and cannot be solved by Z3 or cvc5. (Z3 returns {\tt unknown} as it does not support the {\tt str.replace\_re\_all} function well; cvc5 fails to solve the problem in 300 seconds.) 
%To tackle the challenge, we utilize and extend the framework established in the string solver {\sf OSTRICH}~\cite{atva2020} to reason about the complex as well as non-standard string operations and obtain a solver called \oursolver. In the end, 
Our new solver $\ostrichseq$ solves Listing~\ref{lst:js_example_constraints} in less than 1~second.
 
%\zhilin{stopped here}

%and then joins the all parts consisting of number characters,  and finally returns its numeric representation. Moreover, the function has a precondition and a postcondition. We use regular expressions to express the precondition and postcondition. $version \in$ \verb|[0-9]+([0-9a-z._ \/\-])*| is the precondition, which checks if the input string contains numbers and each number is followed by a non-numeric character. The postcondition is $version \in$ \verb|[0-9]+(\.[0-9]+)*|, which checks if the output string is a float.

%%%%%%%%%%%%%%%%%%%%%%% removed
%%%%%%%%%%%%%%%%%%%%%%% removed
\hide{
To check the path feasibility of the function, we need to solve the constraints involving string sequences. For example, we need to check whether the length of the sequence \texttt{numbers} is equal to 1 or greater than 1, and we also need to model the \texttt{split} operation and the \texttt{join} operation. The state-of-the-art SMT solvers, such as Z3 and cvc5, are not capable of solving these constraints directly. Our solver $\ostrichseq$, however, can handle these constraints effectively. The constraints generated from the example can be expressed in our theory of string sequences, which includes operations like \texttt{seq.len}, \texttt{seq.++}, \texttt{seq.extract}, \texttt{seq.nth}, \texttt{str.split}, and \texttt{seq.join}. The path feasibility constraints of \textbf{path 2} is shown in Listing \ref{lst:js_example_constraints}. The \texttt{precondition} is a regular expression that checks if the input string contains numbers, and the \texttt{postcondition} checks if the output string is a float. Line 2 ensures that the input string \texttt{version0} satisfies the precondition. Line 3 asserts that \texttt{numbers0} is the result of splitting the input string \texttt{version0} by the regular expression \verb|/[a-z._ \/\-]/i|. Line 4 asserts that the length of \texttt{numbers0} is greater than 1, which corresponds to the condition in the JavaScript code. Line 5-9 constructs the final result of the version string, corresponding to the operations of line 8-10 in the JavaScript code \ref{lst:js_example}. Finally, line 10 checks the negation of the postcondition. If the constraints are satisfiable, it means that there is an input string that satisfies the precondition and produces an output string that does not satisfy the postcondition, indicating a potential bug in the code. Our solver $\ostrichseq$ can solve these constraints within 1 second, while Z3 and cvc5 fail to solve them within a reasonable time.
}
%%%%%%%%%%%%%%%%%%%%%%% removed
%%%%%%%%%%%%%%%%%%%%%%% removed

%\vspace{-2mm}
\section{A Theory of String Sequences ($\arraylogic$)}\label{sec:logic}

%!TEX root = main.tex

A \emph{sequence} over the set $X$ is $(a_0, \ldots, a_{n-1})$, where $a_i \in X$ for each $i \in [0,n-1]$, and $n$ is the \emph{length} of the sequence.
%A \emph{string} is a sequence over %$A$ is 
%a (finite) alphabet $\Sigma$. As usual, $\Sigma^*$ denotes the set of strings over $\Sigma$, and the empty sequence/string is denoted by $\varepsilon$.
We shall focus on \emph{string sequences}, namely, sequences over $\Sigma^*$. We use $()$ to denote the empty sequence.
% for some (finite) alphabet $\Sigma$. 

%We consider the following 
\smallskip
\noindent\emph{Operations on string sequences and strings.} %Let $\Sigma$ be an alphabet and 
We assume two string sequences $s_1 = (u_0, \ldots, u_{m-1})$ and  $s_2 = (v_0, \ldots, v_{n-1})$.
\begin{itemize}%[topsep=0pt,partopsep=0pt]
	\item $s_1 \concat s_2$ is the concatenation of $s_1$ and $s_2$, i.e., $(u_0, \ldots, u_{m-1}, v_0, \ldots, v_{n-1})$.

	      % \item $\transducer(u)$, which applies a \emph{functional} transduction $\transducer$ to $u\in \Sigma^*$ and returns  $v$ such that $(u, v) \in \transducer$. \tl{check}
	      %
	\item $\myat(s_1, i)$ is $u_i$ if $i\in [0,m-1]$, and  undefined otherwise.
	      %
	\item $\myjoin_v(s_1)$ joins the elements in $s_1$ with $v$ as the intermediate string, i.e., \polished{$\myjoin_v(s_1) = u_0 v u_1 v u_2 v \cdots v u_{m-1}$}. 
		      %
	\item $s_1[i \rightarrow u]$ replaces $u_i$ in $s_1$ by $u$ if $i\in [0,m-1]$,  and is $s_1$ otherwise,  i.e.
	\begin{center}
	 $[s_1[i \rightarrow u] = 
	\left \{
	\begin{array}{l l}
	(u_0, \ldots, u_{i-1}, u, u_{i+1}, \ldots, u_{m-1}), & \mbox{ if } i\in [0,m-1]; \\
	s_1, & \mbox{ otherwise}.
	\end{array}
	\right.
	$    
	\end{center}
	      %
	\item $\myfilter_e(s_1)$ filters out all elements in $s_1$ that do not match the regular expression $e$, i.e., $\myfilter_e(s_1) = (u_{i_1}, \ldots, u_{i_k})$  such that $0\le i_1 < \cdots < i_k\le m-1$ and  for all $j \in [0,m-1]$, $u_j \in \Lang(e)$ iff $j \in \{i_1, \ldots, i_k\}$. %$u_j \in \Lang(e)$ iff $j \in \{i_1, \ldots, i_k\}$. 
	      %
	      In particular, if none of $u_0, \ldots, u_{m-1}$ are in $\Lang(e)$, then $\myfilter_e(s_1) = ()$. 
	      %\tl{or $\myfilter_e(s_1) = ()$? it should be a sequence!?}\zhilin{yes, use () here}
%\tl{$\varepsilon$ or $(\varepsilon)$?}\zhilin{empty sequence, $\varepsilon$. $(\varepsilon)$ denotes a sequence of length 1, where the element is the empty string ?}
	      %
	\item $\mysplit_e(u)$ splits a string $u$ into a sequence of substrings according to the regular expression $e$, i.e.,  $\mysplit_e(u) = (u'_1, \ldots, u'_k)$, where $u = u'_1 v_1 u'_2 v_2 \cdots u'_{k-1}v_{k-1} u'_k$ such that (1) for each $i \in [k-1]$, $v_i$ is \polished{the leftmost and longest substring of $u'_i v_i u'_{i+1}\cdots v_{k-1}u'_k$} that matches $e$, and (2) $u'_k$ does not contain any substring that matches $e$.
	      % moreover, $u'_k$ does not contain occurrences of $v$.
	      In particular, if $u$ does not contain any substring that matches $e$, then $\mysplit_e(u) = (u)$.

	\item $s_1[i, j]$ is defined as follows: if $i\in [0,m-1]$ and $j \ge 1$, then $s_1[i, j]$ is  the subsequence $(u_i, \ldots, u_{\min(i+j-1,m-1)})$; if $i\in [0,m-1]$ and $j = 0$, then $s_1[i, j]$ is the empty sequence $()$;  if $i\not\in [0,m-1]$ or $j < 0$, then $s_1[i, j]$ is undefined. 
	%\tl{can we try to avoid $\bot$? tbd}

	      % \item $\mymatch_e(u)$ extracts the leftmost occurrence of $e$ in a string $u$ and returns the sequence $(v)$, where $u = u'_1 v u'_2$ such that $v$ is the leftmost occurrence of $e$ in $u$.  
	      %
	\item $\mymatchall_e(u)$ extracts the substrings of $u$ that match the regular expression $e$, 
	%all occurrences of $e$ in the string $u$,
	 i.e., $\mymatchall_e(u) = (v_1, \ldots, v_{k-1})$ such that $u = u'_1 v_1 u'_2 v_2 \cdots u'_{k-1}v_{k-1} u'_k$ and (1) for each $i \in [k-1]$, $v_i$ is \polished{the leftmost and longest substring of $u'_i v_i u'_{i+1}\cdots v_{k-1}u'_k$}, and (2) $u'_k$ has no substring that matches $e$.

	      % \item $\transducer(s_1)$ applies a transduction $\transducer$ to $s_1$ elementwise, i.e., $(\transducer(u_1), \ldots, \transducer(u_m))$.
	      %subsequence, match, matchAll, transduction, 
\end{itemize}

\polished{Note that we choose the longest match semantics in $\mysplit_e(u)$ and $\mymatchall_e(u)$
just for illustrating our approach, where
the shortest match semantics can be captured by adapting the automata construction in the pre-image computation.}

%A string over an alphabet $\Sigma$ is a sequence of symbols from $\Sigma$. 
\smallskip
\noindent\emph{The logic $\arraylogic$ of string sequences.} $\arraylogic$   
has three data types: $\inttype$ (integer), $\strtype$ (strings), and $\seqtype$ (sequences).
%In this paper, $\seqtype$ is seen as a composite data type \tl{not sure about this term} where the elements are of type $\strtype$. 
%
We use $u, v, \ldots$ (resp. $\svaru, \svarv, \ldots$) to denote string constants (resp. variables); $m,n, \ldots$ (resp. $\svarm, \svarn, \ldots$) to denote integer constants (resp. variables);
$s, t, \ldots$ (resp. $\svars,\svart, \ldots$) to denote sequence constants (resp. variables); and $e$ to denote regular expressions.
%We define a theory $\arraylogic$ with three sorts: \textbf{string}, \textbf{integer} and \textbf{array}. The \textit{elements} of \textbf{array} are in \textbf{string} sort, and the $index$ of \textbf{array} is in \textbf{Integer} sort. 
%\begin{figure}[htbp]
%	\renewcommand{\arraystretch}{1.3}

The syntax of $\arraylogic$ is defined as follows:
\[
\begin{tabular}{c l l l r}
		%        $u, v, \cdots$     &  &                                                                          & string constants           \\
		%        $\svarx, \svary, \cdots$     &  &                                                                          & string variables           \\
		%        $m,n, \cdots$     &  &                                                                          & integer constants         \\
		%        $\ivarx, \ivary, \cdots$     &  &                                                                          & integer variables          \\
		%        $s, s', \cdots$     &  &                                                                          & sequence constants            \\
		%        $\seqvara,\seqvarb, \cdots$     & &                                                                          & sequence variables            \\
	$~$ &	$it$      & $ \eqdef $ & $m \mid \svarm \mid it+it \mid it-it \mid \mystrlen(strt)\mid \myseqlen(seqt) $                                                      & integer terms  \\
		&	$strt$    & $ \eqdef $ & $u \mid \svaru \mid strt \concat strt \mid \myat(seqt, it) \mid \myjoin_u(seqt)    $                      & string terms   \\
		&	$seqt$    & $ \eqdef $ & $(strt) \mid s \mid \svars \mid seqt \concat seqt \mid seqt[it \rightarrow strt] \mid  \myfilter_e(seqt) \mid  $ &                \\
			&          &            & $\ \ seqt[it, it]  \mid \mysplit_e(strt) \mid \mymatchall_e(strt) $                              & sequence terms \\
		%        $a$      & $::=$    & $\seqvara = \seqvarb[m: n] \mid A = \transducer(B) \mid x=\myread(A, m) \mid$          &                           \\
		%                  &          & $A = \mywrite(B, x ,m) \mid x = \myjoin(A,u) \mid A = \mysplit(x, e) \mid$ &                           \\
		%                  &          & $x= y\cdot z \mid x = \transducer(y)\mid x\in e$                         & atomic formulas        \\
		&	$\varphi$ & $ \eqdef $ & $ it \iop it  \mid seqt = seqt \mid strt = strt \mid strt \in e  \mid \varphi \wedge \varphi $                   & formulas
	\end{tabular}
    \]
where  $\iop\, \in \{=, \neq, <, \le, >, \ge\}$.
%	\caption{Syntax of $\arraylogic$}
%	\label{fig:syntax}
%\end{figure}
%
%The syntax of $\arraylogic$ is defined by the BNF in Figure~\ref{fig:syntax}, 

Here, $it$, $seqt$ and $strt$ denote integer, sequence and string terms, respectively. In particular,
the integer terms are linear arithmetic expressions where $\myseqlen(\svars)$ (resp. $\mystrlen(\svaru)$) denotes the length of a sequence $\svars$ (resp. string $\svaru$);
the sequence terms are constructed from a singleton sequence, %(sequences of length one), 
sequence constants and sequence variables by applying the aforementioned sequence/string operations.
%: concatenation ($\concat$), write ($seqt[it \rightarrow strt]$, writing a string $strt$ to a position $it$ in the sequence $seqt$), split operation ($\mysplit(strt, u)$),  subsequence ($seqt[it, it]$), $\mymatch(strt, e)$, $\mymatchall(strt, e) $, and $\transducer(seqt)$. 
%
A $\arraylogic$ formula is a conjunction of atomic formulas, each of which is of the form $it_1 \iop it_2$ where $\iop\, \in \{=, \neq, <, \le, >, \ge\}$,  a sequence equality $seqt_1 = seqt_2$, or a string equality $strt_1 = strt_2$.
%
%For each regular expression $e$, the language defined by $e$, denoted by $\Lang(e)$, is defined recursively. For instance, $\Lang(\emptyset) = \emptyset$, $\Lang(\varepsilon) = \{\varepsilon\}$, $\Lang(a) = \{a\}$, and $\Lang(e_1 \cdot e_2) = \Lang(e_1) \cdot \Lang(e_2)$, $\Lang(e_1 + e_2) = \Lang(e_1) \cup \Lang(e_2)$, and so on. 
%
Note that sequence and string inequalities can be expressed as well. %seen as macros, with the help of existential quantifiers. 
For instance (using disjunctions and quantifiers for the sake of presentation):
%
%$$\svarx \neq \svary ~\equiv~ \exists \svarx_1, \svarx_2, \svary_2.\ \bigvee \limits_{a, b \in \Sigma, a \neq b} (\svarx = \svarx_1 a \svarx_2 \wedge \svary = \svarx_1 b \svary_2) \vee \exists \svarz.\ \bigvee_{a \in \Sigma} (\svarx = \svary a \svarz \vee \svary = \svarx a \svarz),$$
%\tl{how about $\varepsilon$ and $a$?}
%\tl{does not read right ?}%
 \begin{align*}
    \svarx \neq \svary & \equiv& \exists \svarx_1, \svarx_2, \svary_2. \bigvee \limits_{a, b \in \Sigma, a \neq b} (\svarx = \svarx_1 a \svarx_2 \wedge \svary = \svarx_1 b \svary_2) \vee \exists \svarz. \bigvee_{a \in \Sigma} (\svarx = \svary a \svarz \vee \svary = \svarx a \svarz) \\
		\seqvara \neq \seqvarb &\equiv & (\exists \seqvara_1, \seqvara_2, \seqvarb_2, \svarz_1, \svarz_2.\ \seqvara = \seqvara_1 \concat (\svarz_1) \concat \seqvara_2 \wedge \seqvarb = \seqvara_1 \concat (\svarz_2) \concat \seqvarb_2 \wedge \svarz_1 \neq \svarz_2)\ \vee \\
		                   &   & %\hspace{4mm} 
                               \exists \seqvarg, \svarz.\ \seqvara  = \seqvarb \concat (\svarz) \concat \seqvarg \vee    \seqvarb  = \seqvara \concat (\svarz) \concat \seqvarg.
\end{\align*}
%
\polished{Though rewriting inequalities produces disjunctions, they can still be handled in the DPLL(T) framework~\cite {dpll_t}. Hence, we focus on the disjunction-free fragment.}

%%%%%%%%%%%%%%%%%%%%%%%%%%%%%%
%\hide{
%	A $\arraylogic$ formula $\varphi$ is a conjunction of formulas of the form $t_1\circ t_2$, $A = B[m:n]$, $A = \transducer(B)$, $x = \myread(A,m)$, $A = \mywrite(B, x, m)$, $x=\myjoin(A,u)$, $A=\mysplit(x,e)$, $x = \transducer(y)$, $x=y\cdot z$ or $x\in e$. Definitely, $t$ is a integer term built from the constant $m$, the varible $i$, the size of array $\mylen(A)$ and the length of string $\mylen(x)$ by operators $+$ and $-$; $\circ \in \{=,\not=, <, \leq, >, \geq \}$ is the set of arithmatic binary predicates on integer terms $t_1$ and $t_2$; $e$ is a regular expression; $B[m:n]$ is the sub-array of $B$ containing elements between $m$ and $n$; $\transducer(B)$ modifies the elements of $B$ by transducer $\transducer$ and returns the result; $\myread(A, m)$ returns the $m$-th element of $A$; $\mywrite(B, x ,m)$ returns the array $B$ after $x$ is written to $m$-th position of $B$; $\myjoin(A,u)$ joins all of the elements in $A$ by string $u$; $\mysplit(x,e)$ splits $x$ into an ordered string array by regex $e$; $\transducer(y)$ returns the output of the transducer $\transducer$ running on $y$; $y\cdot z$ is the concatenation of $y$ and $z$.
%}
%%%%%%%%%%%%%%%%%%%%%%%%%%%%%

\begin{proposition}\label{prop-undec}
	The satisfiability of $\arraylogic$ is undecidable.
\end{proposition}
%Idea: replaceAll where the replacement strings are constants can be captured by transducers. 
%
%From Proposition~\ref{prop-undec}, we know that the satisfiability of $\arraylogic$ is undecidable. In the sequel, we show that for the straight-line fragments of $\arraylogic$, the satisfiability is decidable. 

We define the \emph{straight-line fragment} $\slarraylogic$ of $\arraylogic$. It is easy to see that by introducing fresh %sequence/string 
variables, we can rewrite a $\arraylogic$ formula to a form in which the left-hand side of each equality $seqt_1 = seqt_2$ (resp. $strt_1 = strt_2$) is a sequence (resp. string) variable.
%
%i.e., the string/sequence equalities are in one of the following forms,
%\[
%	\begin{array} {l}
%		\svars = (\svaru) \mid \svars = s \mid \svars = \svart \mid \svars = \svart_1 \concat \svart_2 \mid \svars = \svart[it \rightarrow \svaru] \mid \\
%		  \svars  = \myfilter_e(\svart) \mid \svars = \svart[it_1, it_2] \mid \svars = \mysplit_e(\svaru) \mid \svars = \mymatchall_e(\svaru) \mid \\
%		\svaru = u \mid \svaru = \svarv \mid \svaru = \svarv_1 \concat \svarv_2 \mid \svaru = \myat(\svars, it) \mid \svaru = \myjoin_u(\svars).
%	\end{array}
%\]
Note that the original and rewritten formulas are equisatisfiable. %rewriting preserves the satisfiability. 
Henceforth, we assume that all sequence/string equalities satisfy this constraint unless otherwise stated. %of the aforementioned forms. 

%It is easy to see that for a sequence equality $seqt_1 = seqt_2$, we can introduce a fresh sequence variable $\seqvara$ and rewrite  it as $\seqvara= seqt_1 \wedge \seqvara = seqt_2$. Similarly, for a string equality $strt_1 = strt_2$, we can introduce a fresh string variable $\svarx = strt_1 \wedge \svarx = strt_2$. Note that the rewritings preserve the satisfiability. As a result, we can assume that all the sequence equalities $seqt_1 = seqt_2$ (resp. string equalities $strt_1 = strt_2$) satisfy that $seqt_1$ is a sequence variable (resp. $strt_1$ is a string variable). 

\begin{definition}[Straight-line fragment]\label{def:straight_line}
	Given a $\arraylogic$ formula $\varphi$, 
	 the \emph{dependency graph} $G_\varphi$ of  $\varphi$
	 is a directed graph $(V,E)$ such that
	%	we define the \emph{dependency graph} $G_\varphi$ of  $\varphi$, such that 
	$V$ is  the set of string and sequence variables in $\varphi$, and % the graph comprising 
	 $(v_1, v_2)\in E$ %such that $v_1,v_2$ are sequence or string variables, 
	iff  $v_1 = {\tt rhs}$ is a sequence or string equality in $\varphi$ and $v_2$ occurs in ${\tt rhs}$ except for integer terms.
	%
	%For instance, from the equality $\seqvara = \seqvarb[1, \mylen(\svarx)]$, we know that $(\seqvara, \seqvarb)$ is an edge in $G_\varphi$, while $(\seqvara, \svarx)$ is not. 

$\varphi$ is \emph{straight-line} if $\varphi$ is disjunction-free, and each string/sequence variable occurs as the left-hand side of equalities at most once; moreover, $G_\varphi$ is acyclic.

%	$\varphi$ is \emph{straight-line} if $G_\varphi$ is acyclic.
\end{definition}
%Let us use $\slarraylogic$ to denote the straight-line fragment of $\arraylogic$.

\begin{example}
$\svars = \mysplit_e(\svaru) \wedge \svart = \svars[1, \myseqlen(\svars) - 1]$ is  straight-line, while  $\svars = \svars[1, \myseqlen(\svars) - 1]$ is not, since there is a self-dependency on $\svars$.
\end{example}

%Straight-line fragments: Write the equalities as $\svarx = strt$ or $\seqvara = seqt$. The dependency graph should be acyclic (where the string variables and sequence variables are all taken into consideration). 

%\zhilin{stopped here}

%%%%%%%%%%%%%%%%%%%%%%%%%%%%%%%%%%
%%%%%%%%%%%%%%%%%%%%%%%%%%%%%%%%%%
\hide{
	%\subsection{Definition of Straight-Line Formula }
	\begin{definition}[Relational constraints ]
		Relational constraints $\phi$ are defined by the following rules:
		\begin{align*}
			\phi \ ::= \  & A = B[m:n] \mid A = \transducer(B) \mid x=\myjoin(A,u) \mid A = \mysplit(x,e) \mid                       \\
			              & x=\myread(A, m) \mid A = \mywrite(B, x ,m) \mid x = \transducer(y) \mid x=y\cdot z \mid \phi \wedge \phi
		\end{align*}
	\end{definition}
	\begin{definition}[Straight-line relational constraints]
		A relational constraint $\phi$ is straight-line if $\phi = \bigwedge\limits_{1\leq i \leq m} \chi_i = P_i$ such that
		\begin{itemize}
			\item $\chi_1,\cdots,\chi_m $ are mutally distinct string varibles and array varibles.
			\item For each $i\in [m] $, all the varibles in $P_i$ are in the set $\{\chi_1,\cdots,\chi_{m-1} \}$.
		\end{itemize}
	\end{definition}
	\begin{definition}[Straight-line formula]
		A formula $F$ of $\arraylogic$ is straight-line iff all relational constraints in $F$ are straight-line.
	\end{definition}
}
%%%%%%%%%%%%%%%%%%%%%%%%%%%%%%%%%%
%%%%%%%%%%%%%%%%%%%%%%%%%%%%%%%%%%

%\vspace{-2mm}
\section{The Decision Procedure} \label{sect:procedure}

%!TEX root = main.tex


\begin{theorem}\label{thm-main}
	The satisfiability of $\slarraylogic$ is decidable.
\end{theorem}

The general idea of the decision procedure is to encode each string sequence as a string, based on which all string sequence operations in $\slarraylogic$ can be transformed into string operations. %whose pre-images can be effectively computed. 
It utilizes the framework~\cite{atva2020} for solving straight-line string constraints with integer data type,  \polished{as illustrated in Section~\ref{sect:example}}. There are, however, certain technical challenges. For instance, %We show that while some operations in $\slarraylogic$ 
%operations such as 
$seqt[it \rightarrow strt]$ and $seqt[it, it]$ cannot be captured directly by standard string operations,  
%therein. In the sequel, we introduce several new string operations 
so new ones have to be provided whose pre-images need to be computed. %so that the framework in~\cite{atva2020} can still be applied. %in \cite{atva2020}, the framework therein is general enough and can still be extended to tackle these operations. 
%
%In the sequel, we first show how to encode the string sequences as strings and operations on string sequences as string operations, then show how to compute the pre-images of these operations. 

% \vspace{-3mm}
\subsection{From string sequences to strings}\label{subsect:seq2string}
%\denghang{Stop here}

%!TEX root = main.tex

 
We fix a fresh symbol $\separator\notin\Sigma$ and let $\Sigma':=\Sigma \cup \{\separator\}$.
Each string sequence  $s = (u_1, \cdots, u_m)$ over $\Sigma$ is encoded as a string  $\enc(s):=\separator u_1\separator \cdots \separator u_m \separator$ over $\Sigma'$. % the (extended) alphabet $\Sigma \cup \{\separator\}$.  
Note that the strings %encodings of string sequences are not arbitrary strings from $(\Sigma \cup \{\separator\})^*$, but the strings that match 
resulted from the encoding should match the regular expression $\separator (\Sigma^* \separator)^*$ and the string $\separator$ encodes the empty sequence.
%For a sequence $s$, let $\enc(s)$ denote the string encoding of $s$. 

With this encoding of sequences as strings, each string sequence operation is also transformed into the corresponding \emph{string operation}. 
 %all the sequence operations can be turned into string operations. 
%
\begin{description}
\item[sequence concatenation $s_1 \concat s_2$:] %is turned into the 
string concatenation  $\enc(s_1) \concat \enc(s_2)$. 
%
\item[sequence  write {$s[i \rightarrow u]$}: ] 
%string operation 
$\mywrite(\enc(s), i+1, u)$, which replaces the substring of $\enc(s)$ between the $(i+1)$-th occurrence of $\separator$ and the $(i+2)$-th occurrence of $\separator$, by $u$, if $i\in[0,n-2]$, where $n$ is the number of occurrences of $\separator$ in $\enc(s)$, and is undefined otherwise (i.e. $i\not\in[0,n-2]$). Note that this operation is not a typical string replace operation.  (Recall that the indices of elements in sequences start from $0$ and $n$, i.e. the number of occurrences of $\separator$ in $\enc(s)$, is one more than the length of $s$.) 
%Although this operation is not a typical replace operation, we can show that its pre-image can still be effectively computed. 

\item[sequence filter operation $\myfilter_e(s)$:] %string operation 
$\myfilter_e(\enc(s))$, which removes every substring of $\enc(s)$ between two consecutive occurrences of $\separator$ that does not match the regular expression $e$. (For instance, if $e = a^* b$ and $s = (ab, ac, aab)$, then $\myfilter_e(\enc(s)) = \separator ab \separator aab \separator$.)
%
\item[subsequence operation {$s[i, j]$}:] %$s[i, j]$ is turned into 
$\subseq(\enc(s), i+1, j)$, which keeps only the substring of $\enc(s)$ between the $(i+1)$-th occurrence and the $\min(i+j, n-1)$-th occurrence of $\separator$, if $i\in[0,n-2]$, where $n$ is the number of occurrences of $\separator$ in $\enc(s)$, and  is undefined otherwise (i.e. $i\not\in[0,n-2]$). 
%
\item [$\mysplit_e(u)$:] string operation $\mysplitstr_e(u)$ that transforms $u = u'_1 v_1 u'_2 \cdots u'_{k-1} v_{k-1} u'_k$ into  $\separator u'_1 \separator u'_2 \cdots u'_{k-1} \separator u'_k \separator$, where for each $i \in [k-1]$,  $v_1$ is the leftmost and longest matching of $e$ in $u'_{i} v_i u'_{i+1}$, moreover,  $u'_k$ does not contain any substring that matches $e$. 
%
%(Note that here we abuse the notation a bit by using $\mysplit_e$ to denote both the operation of splitting a string into a sequence and the operation of transducing a string into another string. Similarly for the operation $\mymatchall_e$ in the sequel.)
%
\item[$\mymatchall_e(u)$:] %string operation 
$\mymatchallstr_e(u)$ that transforms $u = u'_1 v_1 u'_2 v_2 \cdots u'_{k-1}v_{k-1} u'_k$ into 
$\separator v_1 \separator \cdots \separator v_{k-1} \separator$, where for each $i \in [k-1]$, $v_i$ is the leftmost and longest occurrence of $e$ in $u'_i v_i u'_{i+1}$, moreover, $u'_k$ does not contain any substring that matches $e$.

%\item The transduction operation $\transducer(\seqvarb)$ is turned into a string operation $\transducer_{\separator}(\enc(\seqvarb))$, where $\transducer_{\separator}$ is the transducer that replaces every substring between  two occurrences of $\separator$, say $u$, by $\transducer(u)$.
%
\item[sequence read operation $\myat(s, i+1)$:] $\elem(\enc(s), i+1)$ that keeps only the substring between the $(i+1)$-th occurrence  and the $(i+2)$-th occurrence of $\separator$, if $i\in[0,n-2]$, where $n$ is the number of occurrences of $\separator$ in $\enc(s)$, and is undefined otherwise (i.e. $i\not\in[0,n-2]$). (For instance, if $s = (ab, ac)$ and $i = 0$, then $\elem(\enc(s), 1) = \elem(\separator ab \separator ac \separator, 1) = ab$.)
%
\item[sequence join operation $\myjoin_u(s)$:] $\myjoin_u(\enc(s))$ that removes the first and the last occurrences of $\separator$ and replaces all the remaining occurrences of $\separator$ by $u$.
\item[sequence length operation $\myseqlen(s)$:] $\myseqlen(\enc(s))$ - 1 where $\myseqlen(\enc(s))$ counts the number of occurrences of $\separator$ in $\enc(s)$.
\end{description}

%After encoding string sequences into strings, 
Hence we obtain atomic string constraints %where the string equalities are in one 
of the following forms, 
\begin{equation*}\label{eqn-str-enc}
\begin{array} {l}
\svarx = u \mid \svarx = \svary \mid \svarx = \svary \concat \svarz \mid \svarx = \transducer(\svary) \mid \svarz = \mywrite(\svarx, it, \svary) \mid  \svary  = \myfilter_e(\svarx)  \mid   \\
\svary = \mysplitstr_e(\svarx) \mid \svary = \subseq(\svarx, it_1, it_2) \mid   
 \svary = \mymatchallstr_e(\svarx) \mid \svary = \elem(\svarx, it)  \mid \\
 \svary = \myjoin_u(\svarx) \mid it = \mystrlen(\svarx) \mid it = \myseqlen(\svarx),
\end{array}
\end{equation*}
where $\transducer$ is a finite-state transducer. 

We denote by $\exstrlogic$ %(where $\sf Ex$ means ``Extended'') to denote 
the class of string constraints which are %defined as 
positive Boolean combinations (no negation) of the atomic string constraints, and  %aforementioned string equalities. Moreover, let us use 
$\slexstrlogic$ denotes its straight-line fragment.  % of $\exstrlogic$, 
%where string inequalities do not occur and there are no cyclic dependencies between string variables. 
For each formula $\varphi$ in $\slarraylogic$, $\enc(\varphi)$ is the resulting formula in $\exstrlogic$. We have the following result.

\begin{proposition} %\tl{this is to be improved}
	% $\arraylogic$ constraints into $\exstrlogic$ preserves the satisfiability. 
For each  constraint $\varphi$ in $\arraylogic$, 
	$\varphi$ and $\enc(\varphi)$ are equisatisfiable.
	Moreover, if $\varphi$ is in $\slarraylogic$, then $\enc(\varphi)$ is in $\slexstrlogic$.
\end{proposition}
%It is not hard to see that the aforementioned encoding of 
\begin{example}
Consider the $\slarraylogic$ constraint $\varphi:= (\svars_1 = \svars_0\cdot(u,v)\wedge\myseqlen(\svars_1) < 2)$ where $\svars_0$ is a string variable and $u,v$ are string constants. 
Then, its string encoding is $\enc(\varphi) := (\enc(\svars_1) = \enc(\svars_0)\cdot \separator u \separator v \separator \wedge \myseqlen(\enc(\svars_1)) - 1  < 2)$, which is unsatisfiable, since there are at least three occurrences of $\separator$ in $\enc(\svars_1)$. \qed
\end{example}

The satisfiability of $\slarraylogic$ is now reduced to that of $\slexstrlogic$. %Our goal is then to apply the framework in~\cite{atva2020} to solve the $\slexstrlogic$ constraints. At first, 
It is easy to observe that the operations $\myfilter_e$, $\mysplitstr_e$, $\mymatchallstr_e$ and $\myjoin_u$ can all be represented by a finite-state transducer, %seen as special cases of transducers $\transducer$, thus their pre-images can be effectively computed according to the 
so the results in~\cite{atva2020} are directly applicable. 
In the next section, %It remains to 
we show that the pre-images of the remaining string operations, i.e., %that are not captured directly by the operations in~\cite{atva2020}, that is, 
$\mywrite$, $\subseq$, $\elem$ and $\myseqlen$, 
%satisfy that their  
can be effectively computed. 




% \vspace{-2mm}
\subsection{Computing the pre-images under $\mywrite$, $\subseq$, $\elem$ and $\myseqlen$} \label{subsect:preimage}

%!TEX root = main.tex

%In this section, we show how to compute the pre-images of the string operations in Equation~(\ref{eqn-str-enc}). The operations $\myfilter_e$,  $\mysplit_u$, $\mymatch_e$, $\mymatchall_e$ and $\myjoin_u$ can be easily encoded in transducers, the computations of their pre-images can be derived by the work \cite{atva2020}. (The details can be found in Appendix~\ref{???}.) We focus on $\mywrite$, $\subseq$ and $\elem$. 

\cite{atva2020} introduced cost-enriched finite automata (CEFA for short) and showed that the pre-images of all the considered string operations involving the string and integer data types only can be effectively computed. 
% 
%The decision procedure in~\cite{atva2020} introduced cost-enriched finite automata (CEFA for short) and show that the pre-images of all the string operations involving the string and integer data type can be effectively computed.
The OSTRICH %string constraint solving 
framework 
%If all the string operations in a 
essentially informs that once the pre-image is computable, the straight-line %string constraint satisfy this property, then the constraint 
fragment can be solved by propagating the constraints encoded in CEFA  %of the form $\svaru \in \aut$ 
backwards. Our decision procedure for  $\slexstrlogic$ constraints follow this approach. 

We first recall the definition of CEFAs.
%Recall the cost-enriched finite automata (CEFAs~\cite{atva2020}). 
Intuitively, CEFAs add write-only cost registers to finite state automata, where ``write-only'' means that the cost registers can only be written/updated but cannot be read, i.e., they cannot be used in the guards of the transitions. 
%Moreover, an LIA \tl{remember to introduce this earlier} formula on the values of the registers is introduced as an accepting condition. 
%
%Recall that we assume that $\Sigma$ is the alphabet of the strings/elements in the sequences and $\separator \not \in \Sigma$ is a separator used to encode string sequences as strings. 


%\begin{definition}[Cost-Enriched Finite Automaton]
A CEFA $\aut$ over $\Sigma'$ is a  tuple $(R, Q, \Sigma', \delta, I, F)$ where
	\begin{itemize}
		\item $R = \{r_1, \cdots, r_k\}$ is a finite set of registers,
		\item $Q, I, F$ are the same as in {\nfa}, i.e., the set of states, the set of initial states, and the set of final states, respectively, 
		%    \item $R = (r_1\cdots r_n)$ is a vector of mutually distinct cost registers,
		\item $\delta \subseteq Q \times \Sigma' \times Q \times \Int^R$ is a transition relation, where $\Int^R$ denotes the updates on the values of registers. (Recall that $\Sigma':=\Sigma \cup \{\separator\}$ is the extended alphabet.)
		%  set whose elements are tuples $(q, c, q', \myvec{v})$ where $q, q'$ are states of $Q$, $c$ is a letter in alphabet $\Sigma\cup\{\epsilon\}$ and $\myvec{v}$ is the cost update function for registers, which is an integer vector whose $i$-th element is the incremental value of register $r_i$.  $(q, a, q', \myvec{v})$ is written as $q\xrightarrow[\myvec{v}]{a} q'$ for readability.
%		\item $acc \in \Phi(R)$ is an LIA formula specifying an accepting condition.
		%    a linear integer arithmetic constraint function on final states. $\theta$ is called \emph{accepting condition}.
	\end{itemize}
%\end{definition}
We write $R_\aut$ for the set of registers of $\aut$. %We assume a linear order on $R$ and write 
Typically, $R_\aut$ is represented as a vector $(r_1, \cdots, r_k)$. Accordingly, we write assignments $r_i:= r_i + v_i$ for all $i\in [k]$ %element of $\Int^R$ 
simply as 
%its values are represented as 
$\vec{v}=(v_1, \cdots, v_k)$, i.e., $r_i$ is to be incremented by $v_i$ for each $i \in [k]$. We also write a transition $(q, a, q', \vec{v}) \in \delta$ as $q \xrightarrow[\vec{v}]{a} q'$. In particular, $(q,a,q')$ stands for the transition $(q,a,q',())$ without any updates.
%The definition of CEFA above is slightly different from the definition in~\cite{atva2020} in the sense that accepting conditions on register values
%are attached to final states. 

%The semantics of CEFA is defined as follows. 
Let $\aut = (R, Q, \Sigma', \delta, I, F)$ be a CEFA. 
%A \emph{configuration} of $\aut$ is a pair $(q, \vec{v})$ where $q \in Q$ and $\vec{v}$ is a vector denoting the values of registers.  
%An \emph{initial configuration} of $\aut$ is $(q_0, \vec{0})$ with $q_0 \in I$, where the value of each register is zero. 
A \emph{run} of $\aut$ on a string $w = a_1 \cdots a_n$ over $\Sigma'$ is a sequence $q_0 \xrightarrow[\myvec{v_1}]{a_1} q_1 \cdots q_{n-1}\xrightarrow[\myvec{v_n}]{a_n} q_n$ such that $q_0 \in I$ and $q_{i-1} \xrightarrow[\myvec{v_i}]{a_i} q_i$ for each $i \in [n]$. 
%
%The run $q_0 \xrightarrow[\myvec{v_1}]{a_1} q_1 \cdots q_{n-1}\xrightarrow[\myvec{v_n}]{a_n} q_n$ 
It is \emph{accepting} if $q_n \in F$, and 
%
the vector $\myvec{c}= \sum_{j \in [n]} \myvec{v_j}$ is defined as the \emph{cost} of an accepting run. %$q_0 \xrightarrow[\myvec{v_1}]{a_1} q_1 \cdots q_{n-1}\xrightarrow[\myvec{v_n}]{a_n} q_n$. 
(Note that all registers are initialized to zero.) 
%
%and updated to $\sum \limits_{j \in [n]} \myvec{v_j}$ after all the transitions in the run are executed. 
We write $\myvec{c} \in \aut(w)$ if there is an accepting run of $\aut$ on $w$ whose cost is $\myvec{c}$.  

The semantics of a CEFA $\aut$, denoted by $\Lang(\aut)$, is defined as the set of pairs $\{(w, \myvec{c}) \mid  \myvec{c} \in \aut(w)\}$.
%
In particular, if $I \cap F \neq \emptyset$, then $(\varepsilon, \myvec{0}) \in \Lang(\aut)$. Moreover, the \emph{output} of a CEFA $\aut$, denoted by $\cefaout(\aut)$, is defined as $\{\myvec{c} \mid  \exists w.\ \myvec{c} \in \aut(w)\}$.

%%%%%%%%%%%%%%%%%%%%%%%%%%%%%%%%%%%%%%%%%%%%%%%%%%%%%%%%%%%%%%%%%%%%%%%%%%%%%%%%%
The following example illustrates the backward propagation. 

\vspace*{-2mm}
%\begin{example}
%Consider the equality $\svarv = f(\svaru_1, it_1, \svaru_2, it_2)$ (where $\svaru_1,\svaru_2$ are string variables and $it_1, it_2$ are integer terms) and the CEFA $\aut = (R, Q, \Sigma', \delta, I, F)$, where $R= \{r_1\}$. %(that is, $R$ contains only one register). 
%Suppose that $f^{-1}(\aut)$ is a finite collection of tuples $(\autb_1, \autb_2, t_1)$, where $\autb_1$ and $\autb_2$ are CEFAs such that $R_{\autb_1}=(r'_{1,1}, r^{(1)}_1)$ and $R_{\autb_2} = (r'_{2,1}, r^{(2)}_1)$ and $t_1$ is an LIA formula over $r^{(1)}_1$ and $r^{(2)}_1$. Intuitively, $r'_{1,1}$ and $r'_{2,1}$ are the fresh registers corresponding to the two integer parameters,  and $r^{(1)}_1, r^{(2)}_1$ are the fresh registers introduced in $\autb_1$ and $\autb_2$ respectively for restoring $r_1$. Then the backward propagation of the constraint $\svarv \in \aut$ 
%%with respect to the equality $\svarv = f(\svaru_1, it_1, \svaru_2, it_2)$ 
%wrt. $f$ 
%nondeterministically chooses a tuple $(\autb_1, \autb_2, t_1)$, removes $\svarv = f(\svaru_1, it_1, \svaru_2, it_2)$ from the string constraint, and adds $t_1 \wedge it_1 = r'_{1,1} \wedge it_2 = r'_{2,1}$ to it. \qed
%\end{example}
\begin{example}
	Consider the equality $\svarv = f(\svaru_1, it_1, \svaru_2, it_2)$ (where $\svaru_1,\svaru_2$ are string variables, $it_1, it_2$ are integer terms, and $f$ is a function) and the CEFA $\aut = (R, Q, \Sigma', \delta, I, F)$, where $R= \{r_1\}$. %(that is, $R$ contains only one register). 
	%Moreover, let $t = r^{(1)}_1 + r^{(2)}_1$ which intuitively describes how $r_1$ can be calculated from $r^{(1)}_1$ and $r^{(2)}_1$. 
	Then an $R$-cost-enriched pre-image of $\aut$ under $f$, denoted by $f^{-1}_{R}(\aut)$, is a pair $(\cerl, t)$ where $\cerl \subseteq (\Sigma')^\ast \times (\intnum \times \intnum) \times (\Sigma')^\ast \times (\intnum \times \intnum)$, $t$ is an LIA term over $\{r^{(1)}_1, r^{(2)}_1\}$, and $\cerl$ satisfies that for every $(v, n') \in \Lang(\aut)$, there is $(u_1, n_1, n'_1, u_2, n_2, n'_2) \in \cerl$ such that $f(u_1, n_1, u_2, n_2) = v$, and  $n'= t[n'_1/r^{(1)}_1, n'_2/r^{(2)}_1]$. 
	%
	We say that a pre-image $f^{-1}_{R}(\aut) = (\cerl, t)$ is CEFA-definable if $\cerl$ can be expressed as a finite collection of CEFA tuples $(\autb_{j, 1}, \autb_{j, 2})$ with $j \in [k]$ such that $R_{\autb_{j,1}}=(r'_{1,1}, r^{(1)}_1)$, $R_{\autb_{j, 2}} = (r'_{2,1}, r^{(2)}_1)$, and $\cerl = \bigcup_{j \in [k]} \Lang(\autb_{j, 1}) \times \Lang(\autb_{j, 2})$.  
	%Suppose $t = r^{(1)}_1 + r^{(2)}_1$ and $f^{-1}_{R, t}(\aut)$ is computed as a finite collection of tuples $(\autb_1, \autb_2)$, where $\autb_1$ and $\autb_2$ are CEFAs such that $R_{\autb_1}=(r'_{1,1}, r^{(1)}_1)$ and $R_{\autb_2} = (r'_{2,1}, r^{(2)}_1)$ and $t_1$ is an LIA formula over $r^{(1)}_1$ and $r^{(2)}_1$. Intuitively, 
	Intuitively, $r'_{1,1}$ and $r'_{2,1}$ are the fresh registers corresponding to the two integer parameters, and $r^{(1)}_1, r^{(2)}_1$ are the fresh registers introduced in $\autb_{j, 1}$ and $\autb_{j, 2}$ respectively for computing $r_1$ according to $t$. Then the backward propagation of the constraint $\svarv \in \aut$ 
	%with respect to the equality $\svarv = f(\svaru_1, it_1, \svaru_2, it_2)$ 
	under $f$ nondeterministically chooses a tuple $(\autb_{j, 1}, \autb_{j, 2})$, removes $\svarv = f(\svaru_1, it_1, \svaru_2, it_2)$ from the string constraint, and adds $\svaru_1 \in \autb_{j, 1} \wedge \svaru_2 \in \autb_{j, 2} \wedge r_1 = t \wedge it_1 = r'_{1,1} \wedge it_2 = r'_{2,1}$ to the string constraint.
	%removes $\svarv = f(\svaru_1, it_1, \svaru_2, it_2)$ from the string constraint, and adds $r_1 = t \wedge it_1 = r'_{1,1} \wedge it_2 = r'_{2,1}$ to the string constraint. \qed
\end{example}
\vspace*{-2mm}



%%%%%%%%%%%%%%%%%%%%%%%%%%%%%%%%%%%%%%%%%%%%%%%%%%%%%%%%%%%%%%%%%%%%%%%%%%%%%%%%%%%%%%%%%%%%%%%%%%%%%%% 

We now show how to compute the pre-images of $\mywrite$, $\subseq$, $\elem$ and $\myseqlen$. Note that the pre-image computation utilizes an NFA $\aut_0$ (resp. $\aut_1$) that formats the strings that represent the string sequences (resp. strings), namely, $\aut_0$ (resp. $\aut_1$) recognizes the language $\separator (\Sigma \separator)^*$ (resp. $\Sigma^*$). 

\paragraph*{Pre-image of $\mywrite$.}
Let $\aut = (R, Q, \Sigma', \delta, I, F)$ be a CEFA. 
Intuitively, $\mywrite^{-1}(\aut)$ is computed as the collection $((\autb_{(p, q)} \cap \aut_0, \aut_{R_2/R}[p, q] \cap \aut_1), t_{(p,q)})_{(p, q) \in Q \times Q}$, where  $\autb_{(p, q)}$, $ \aut_{R_2/R}[p, q]$ and $t_{(p,q)}$ are constructed as follows. 
\begin{itemize}
\item $\autb_{(p, q)}$ is constructed by the following procedure. 
\begin{itemize}
%\item We compute an existential LIA formula $\varphi_{p, q}(\myvec{x}, \myvec{y})$ to capture the reachability relation of $\aut[p, q]$. 
%
\item We create two copies of $R$, say $R_1, R_2$. Moreover, we introduce a fresh cost register, say $r'$, to count the number of occurrences of $\separator$. 
%
\item We run $\aut$ on the input string and increase the cost register $r'$ each time when $\separator$ is read. Moreover, the registers in $R_1$, instead of $R$, are updated in the transitions. 
%
\item When the current symbol is $\separator$, we nondeterministically choose to stop increasing $r'$ and pause the running of $\aut$ at the state $p$. Let us call this position as the pause position.
%
\item Finally, when reading the first $\separator$ after the pause position, we resume the run of $\aut$ at $q$, where the registers in $R_1$ (but not $r'$) are updated. 
%
\end{itemize}

\item $\aut_{R_2/R}[p, q]$ is obtained from $\aut[p, q]$ by replacing each register in $R$ with the corresponding copy in $R_2$, where $\aut[p, q]$ is the sub-automaton of $\aut$ that accepts the strings starting from the state $p$ and ending at the state $q$.
%
\item $t_{(p, q)} := (t_{(p, q), r})_{r \in R}$ and 
$t_{(p, q), r} = r^{(1)} + r^{(2)}$ for each $r \in R$.
%\item $t_{(p,q)} := \bigwedge_{r\in R} r = r^{(1)} + r^{(2)}$, where $r^{(1)}$ and $r^{(2)}$ are the copies of $r$ in $R_1$ and $R_2$ respectively.
\end{itemize}
Note that in the backward propagation of $\aut$ with respect to an equality $\svarz = \mywrite(\svarx, it, \svary)$, the constraint $r' = it $ is added to the LIA constraint to assert that $r'$ is equal to the index $it$. 
%
Formally, $\mywrite^{-1}(\aut)$ is a finite collection of $(\autb_{(p,q)}  \cap \aut_0, \aut_{R_2/R}[p, q] \cap \aut_1, t_{(p,q)})$, where $(p, q) \in Q \times Q$, and $\autb_{(p,q)} = (R_1 \cup \{r'\}, Q', \Sigma', \delta', I', F')$ such that  $Q' = Q \times \{pre, idle, post\}$, $I'  = I \times \{pre\}$, $F'  = F \times \{post\}$, and $\delta'$ comprises, 
\begin{itemize}
\item $((q_1, pre), a, (q_2, pre), (\myvec{v}, 0))$ such that $(q_1, a, q_2, \myvec{v}) \in \delta$ and $a \in \Sigma$ (thus $a \neq \separator$), 
%
\item $((q_1, pre), \separator, (q_2, pre), (\myvec{v}, 1))$ such that $(q_1, \separator, q_2, \myvec{v}) \in \delta$, 
%
\item $((q_1, pre), \separator, (p, idle), (\myvec{v}, 1))$ such that $(q_1, \separator, p, \myvec{v}) \in \delta$,
%
%\item $((q_1, idle), a, (q_1, idle), (\myvec{0}, 0))$ such that $a \in \Sigma$, 
%
%\item $((q_1, idle), \separator, (q_2, post), (\myvec{v}, 0))$ such that $(q, \separator, q_2, \myvec{v}) \in \delta$,
%\item $((q_1, idle), \separator, (q_2, post), (\myvec{v}, 0))$ such that $(q, \separator, q_2, \myvec{v}) \in \delta$,
%
\item $((p, idle), a, (p, idle), (\myvec{0}, 0))$ such that $a \in \Sigma$, 
\item $((p, idle), \separator, (q_2, post), (\myvec{v}, 0))$ such that $(q, \separator, q_2, \myvec{v}) \in \delta$,
%
\item $((q_1, post), a, (q_2, post), (\myvec{v}, 0))$ such that $(q_1, a, q_2, \myvec{v}) \in \delta$ and $a \in \Sigma$,
%
\item $((q_1, post), \separator, (q_2, post), (\myvec{v}, 0))$ such that $(q_1, \separator, q_2, \myvec{v}) \in \delta$.
\end{itemize}
From the construction of $\autb_{(p,q)}$, we can observe that the pair $(p,q)$ should satisfy that in $\aut$, there is a $\separator$-transition into $p$ and a $\separator$-transition out of $q$. Moreover, $q$ should be reachable from $p$ according to the  construction of $\aut_{R_2/R}[p, q] \cap \aut_1$. 
% \textcolor{blue}{Did not see how `` $q$ should be reachable from $p$'' is ensured.}
%Moreover, $acc' = acc$. 

\vspace*{-2mm}
\begin{example}
Let us consider the constraint $\svarv = \mywrite(\svaru, 1, \svary)\ \wedge\ \svarv \in \separator (a^+ \separator)^*\ \wedge\ \mystrlen(\svarv) \ge 4$. Let $\aut = (\{r_1\}, \{q_0, q_1, q_2\}, \Sigma', \delta, \{q_0\}, \{q_1\})$ be the CEFA, where the register $r_1$ is used to record the string length and $\delta$ comprises the transitions $q_0 \xrightarrow[(1)]{\separator} q_1 \xrightarrow[(1)]{a}q_2 \xrightarrow[(1)]{a} q_2 \xrightarrow[(1)]{\separator} q_1$. 
Let us consider the pairs $(p, q)$ in $\aut$ satisfying that there is a $\separator$-transition into $p$ and a $\separator$-transition out of $q$, moreover, $q$ should be reachable from $p$.
It is not hard to observe that there is exactly one such pair, that is, $(q_1, q_2)$. 
Therefore, $\mywrite^{-1}(\aut)$ comprises exactly one tuple $(\autb_{(q_1, q_2)}\cap \aut_0, \aut_{R_2/R}[q_1,q_2] \cap \aut_1, (r^{(1)}_1+r^{(2)}_1))$ such that
%and $(\autb_{(q_2, q_2)}\cap \aut_0, \aut_{R_2/R}[q_2,q_2] \cap \aut_1, (r^{(1)}_1+r^{(2)}_1))$ such that  
%\begin{itemize}
%\item 
$\autb_{(q_1,q_2)} = (\{r^{(1)}_1, r' \}, Q', \Sigma', \delta', I', F')$ where 
\begin{itemize}
\item $Q' = \{q_0, q_1, q_2\} \times \{pre, idle, post\}$, $I' = \{(q_0, pre)\}$, $F' = \{(q_1, post)\}$, and 
\item $\delta' $ comprises the following transitions 
\begin{itemize}
\item $((q_1, pre), a, (q_2, pre), (1,0))$, $((q_2, pre), a, (q_2, pre), (1,0))$, 
%$$(q_0, pre)\xrightarrow[(1,1)]{\separator}(q_1, pre)\xrightarrow[(1, 0)]{a}(q_2, pre) \xrightarrow[(1, 0)]{a}(q_2, pre) \xrightarrow[(1, 1)]{\separator}(q_1, pre),$$ 

\item $((q_0, pre), \separator, (q_1, pre), (1,1))$, $((q_2, pre), \separator, (q_1, pre), (1,1))$, 

\item $((q_0, pre), \separator, (q_1, idle), (1, 1))$ (since $(q_0, \separator, q_1, 1) \in \delta$),
%
\item $((q_1, idle), a', (q_1, idle), (0, 0))$ such that $a' \in \Sigma$, 
%
%\item $((q_1, idle), \separator, (q_2, post), (\myvec{v}, 0))$ such that $(q, \separator, q_2, \myvec{v}) \in \delta$,
\item $((q_1, idle), \separator, (q_1, post), (1, 0))$  (since $(q_2, \separator, q_1, 1) \in \delta$),
%
\item $((q_1, post), a, (q_2, post), (1, 0))$ and $((q_2, post), a, (q_2, post), (1, 0))$,
%
\item $((q_2, post), \separator, (q_1, post), (1, 0))$.
%
%$$
%\begin{array}{l}
%(q_0, pre)\xrightarrow[(1,1)]{\separator}(q_1, idle)\xrightarrow[(0, 0)]{\Sigma}(q_1, idle)\xrightarrow[(1, 0)]{\separator}(q_1, post)\\ 
%\xrightarrow[(1, 0)]{a} (q_2, post)  \xrightarrow[(1, 0)]{a} (q_2, post) \xrightarrow[(1, 0)]{\separator} (q_1, post),
%\end{array}
%$$ 
%where $(q_1, idle)\xrightarrow[(1, 0)]{\separator}(q_1, post) \in \delta'$ because $q_2 \xrightarrow[(1)]{\separator} q_1 \in \delta$.   
\end{itemize}
\end{itemize}
To illustrate how $\autb_{(q_1, q_2)}$ and $\aut_{R_2/R}[q_1,q_2]$ work, we can see that $(\separator aa \separator a \separator, 6)$ is accepted by $\aut$, moreover,  for any $b\in\Sigma$,
\begin{itemize}
\item $(\separator ab \separator a \separator, 4, 1)$ is accepted by $\autb_{(q_1, q_2)}$, witnessed by the following sequence of transitions, 
\[
\begin{array}{l}
(q_0, pre) \xrightarrow[(1, 1)]{\separator} (q_1, idle) \xrightarrow[(0, 0)]{a} (q_1, idle) \xrightarrow[(0, 0)]{b} (q_1, idle) \\
\xrightarrow[(1, 0)]{\separator} (q_1, post) \xrightarrow[(1, 0)]{a} (q_2, post) \xrightarrow[(1, 0)]{\separator} (q_1, post),
\end{array}
\]
%
\item and $(aa, 2)$ is accepted by $\aut_{R_2/R}[q_1,q_2]$, witnessed by the following sequence of transitions, 
$
q_1 \xrightarrow[(1)]{a} q_2  \xrightarrow[(1)]{a} q_2
$.  
\end{itemize}
Thus, $\svaru$ and $\svary$  in $\svarv = \mywrite(\svaru, 1, \svary)$
can be $\separator ab \separator a \separator$ and $aa$, respectively.\qed

%
%\item $\aut_{R_2/R}[q_1,q_2]$, 
%\item $\autb_{(q_1, q_1)}$ is similar to $\autb_{(q_1,q_2)}$ and is obtained from $\autb_{(q_1,q_2)}$ by replacing the transition $(q_1, idle)\xrightarrow[(0, 0)]{\separator}(q_2, post)$ in $\delta'$ with $(q_1, idle) \xrightarrow[(0, 0)]{\separator}(q_2, idle)$.
%
%$\autb_{(q_1,q_1)} = (\{r^{(1)}_1, r' \}, Q', \Sigma', \delta',I', F')$ where $Q' = \{q_0, q_1, q_2\} \times \{pre, idle, post\}$, $I' = \{(q_0, pre)\}$, $F' = \{(q_1, post)\}$, and $\delta' $ comprises the following transitions 
%$$(q_0, pre)\xrightarrow[(1,1)]{\separator}(q_1, pre)\xrightarrow[(1, 0)]{a}(q_2, pre) \xrightarrow[(1, 0)]{a}(q_2, pre) \xrightarrow[(1, 1)]{\separator}(q_1, pre),$$ 
%$$
%\begin{array}{l}
%(q_0, pre)\xrightarrow[(1,1)]{\separator}(q_1, idle)\xrightarrow[(0, 0)]{\Sigma}(q_1, idle)\xrightarrow[(0, 0)]{\separator}(q_1, post)\\ 
%\xrightarrow[(1, 0)]{a} (q_2, post) \xrightarrow[(1, 0)]{a} (q_2, post) \xrightarrow[(1, 0)]{\separator} (q_1, post),
%\end{array}
%$$ 
%One tuple of $\mywrite^{-1}(\aut)$ is $(\autb_{(q_1,q_2)}\cap \aut_0, \aut_{R_2/R}[q_1,q_2] \cap \aut_1, true)$ where $\autb_{(q_1,q_2)}$ consist of transitions $(q_0,pre)\xrightarrow[(1)]{\separator}(q_1,pre)\xrightarrow[(0)]{a}(q_2,pre)\xrightarrow[(1)]{\separator}(q_3,pre)$ and $(q_0,pre)\xrightarrow[(1)]{\separator}(q_1, idle)\xrightarrow[(0)]{\Sigma}(q_1,idle)\xrightarrow[(0)]{\separator}(q_2,post)$.
%$\aut_{R_2/R}[q_1,q_2]$ is $q_1\xrightarrow{a}q_2$. After back propagation, we remove the constraint $\svarv = \mywrite(\svaru, 1, \svary)$ and add the constraints $r' = 1\wedge \svaru \in \autb_{(q_1,q_2)}\cap \aut_0\wedge \svary \in \aut_{R_2/R}[q_1,q_2] \cap \aut_1$, where $r'$ is the fresh register introduced in $\autb_{(q_1,q_2)}$. Solving the new constraints, we obtain a solution $\svaru = \separator a \separator$ and $\svary = a$.
\end{example}
\vspace*{-2mm}

%$\myfilter_e$: seen as a special case of transducers

\paragraph*{Pre-image of $\subseq$.}
%$\subseq$: from $\aut$, introduce two fresh cost registers $r_1, r_2$, increment $r_1, r_2$ each time when a $\separator$ is read, nondeterministically choose to stop incrementing $r_1$, start running $\aut$ until nondeterministically choosing to stop incrementing $r_2$, moreover, check whether $\aut$ accepts when stopping incrementing $r_2$. 
%
Let $\aut = (R, Q, \Sigma', \delta, I, F)$ be a CEFA. 
Then $\subseq^{-1}(\aut)$ is computed as $(\autb\cap \aut_0, t)$ such that $t := true$ and $\autb$ introduces two fresh registers $r'_1, r'_2$ to count the two numbers of occurrences of $\separator$ that correspond to the starting position and the length of the subsequence respectively. Moreover, $\autb$ simulates the run of $\aut$ on the substring representing the subsequence returned by $\subseq$. 
%
Formally, $\subseq^{-1}(\aut)$ is computed as $(\autb \cap \aut_0, true)$, where $\autb = (R \cup \{r'_1, r'_2\}, Q \cup \{pre, post\}, \Sigma', \delta', \{pre\}, \{post\})$ such that $pre, post$ are two fresh states denoting the fact that $\autb$ is reading a symbol before and after the subsequence respectively and $\delta'$ comprises the following tuples, 
\begin{itemize}
\item $(pre, a, pre, (\myvec{0}, 0,0))$ such that $a \in \Sigma$ (thus $a \neq \separator$),  
%
\item $(pre, \separator, pre, (\myvec{0}, 1,0))$,
%
\item $(pre, \separator, q, (\myvec{v}, 1,0))$ such that $(p, \separator, q, \myvec{v}) \in \delta$ for some $p \in I$, 
%
\item $(q, a, q', (\myvec{v}, 0, 0))$ such that $(q, a, q', \myvec{v}) \in \delta$ and $a \in \Sigma$, 
%
\item $(q, \separator, q', (\myvec{v}, 0,1))$ such that $(q, \separator, q', \myvec{v}) \in \delta$,
%
\item $(q, \separator, post, (\myvec{v}, 0,1))$ such that $(q, \separator, q', \myvec{v}) \in \delta$ for some $q' \in F$,
%
\item $(post, a, post, (\myvec{0}, 0,0))$ and $(post, \separator, post, (\myvec{0}, 0,0))$ where $a \in \Sigma$.
\end{itemize}
Note that in the backward propagation of $\aut$ with respect to an equality $\svarz = \subseq(\svarx, it_1, it_2)$, the constraint $r'_1 = it_1 \wedge r'_2 = it_2$ is added to the LIA constraint to assert that the value of $r'_1$ and $r'_2$ are equal to the beginning index and the length of the subsequence respectively.

\vspace*{-2mm}
\begin{example}
Let $\svarv = \subseq(\svaru, 2, 1) \wedge \svarv \in \separator (a^+ \separator)^* \wedge \mystrlen(\svarv) \ge 4$ and $\aut = (\{r_1\}, \{q_0, q_1, q_2\}, \Sigma', \delta, \{q_0\}, \{q_1\})$ be the CEFA, where the register $r_1$ is used to record the string length and $\delta$ comprises the transitions $q_0 \xrightarrow[(1)]{\separator} q_1 \xrightarrow[(1)]{a}q_2 \xrightarrow[(1)]{a} q_2 \xrightarrow[(1)]{\separator} q_1$. 
Then $\subseq^{-1}(\aut)$ is $(\autb \cap \aut_0, true)$ such that 
$$\autb = (\{r^{(1)}_1, r'_1, r'_2\}, \{q_0, q_1, q_2, pre, post\}, \Sigma', \delta', \{pre\}, \{post\})$$
where 
$\delta' $ comprises the following transitions 
\begin{itemize}
\item $(pre, a', pre, (0, 0,0))$ such that $a' \in \Sigma$, 
%
\item $(pre, \separator, pre, (0, 1,0))$,
%
\item $(pre, \separator, q_1, (1, 1,0))$ (since $(q_0, \separator, q_1, 1) \in \delta$), 
%
\item $(q_1, a, q_2, (1, 0, 0)$, $(q_2, a, q_2, (1, 0, 0))$, 
%\item $(q, a, q', (\myvec{v}, 0, 0))$ such that $(q, a, q', \myvec{v}) \in \delta$ and $a \in \Sigma$, 
%
\item $(q_2, \separator, q_1, (1, 0, 1))$, 
%$(q, \separator, q', (\myvec{v}, 0,1))$ such that $(q, \separator, q', \myvec{v}) \in \delta$,
%
\item $(q_2, \separator, post, (1, 0, 1))$ (since $(q_2, \separator, q_1, 1) \in \delta$ and $q_1$ is an accepting state in $\aut$), 
%$(q, \separator, post, (\myvec{v}, 0,1))$ such that $(q, \separator, q', \myvec{v}) \in \delta$ for some $q' \in F$,
%
\item $(post, a', post, (0, 0,0))$ and $(post, \separator, post, (0, 0,0))$ where $a' \in \Sigma$.
\end{itemize}
To illustrate how $\autb$ works, we can see that $(\separator a  \separator a a \separator, 4, 2, 1)$ is accepted by $\autb$, witnessed by the following sequence of transitions, 

$
pre \xrightarrow[(0, 1, 0)]{\separator} pre \xrightarrow[(0, 0, 0)]{a} pre \xrightarrow[(1, 1, 0)]{\separator} q_1 \xrightarrow[(1, 0, 0)]{a} q_2 \xrightarrow[(1, 0, 0)]{a} q_2  \xrightarrow[(1, 0, 1)]{\separator} post.
$\qed
%where $\autb$ is 
%$pre\xrightarrow[(1, 0)]{\separator}pre\xrightarrow[(0, 0)]{\Sigma}pre\xrightarrow[(1, 0)]{\separator}q_1\xrightarrow[(0, 0)]{a}q_2\xrightarrow[(0, 1)]{\separator}post\xrightarrow[(0, 0)]{\Sigma'}post$. 
%After back propagation, we remove the constraint $\svarv = \subseq(\svaru, 1, 1)$ and add the constraints $r'_1 = 2\wedge r'_2 = 3\wedge \svaru \in \autb \cap \aut_0$, where $r'_1$ and $r'_2$ are the fresh registers introduced in $\autb$. Solving the new constraints, we obtain a solution $\svaru = \separator\separator a \separator$ which represents the sequence $(\epsilon, a)$.
\end{example}
\vspace*{-4mm}

%$\mysplit_u$: seen as a special case of transducers

%$\mymatch_e$ and $\mymatchall_e$: seen as special cases of transducers


\paragraph*{Pre-image of $\elem$.}
%$\elem$: seen as special cases of transducers
Let $\aut = (R, Q, \Sigma, \delta, I, F)$ be a CEFA. 
The computation of $\elem^{-1}(\aut)$ is similar to that of $\subseq^{-1}(\aut)$. Nevertheless, the construction is slightly different since the output of $\elem$ is a string, instead of a sequence. Intuitively, $\elem^{-1}(\aut)$ is computed as $(\autb \cap \aut_0,t)$ where $t := true$, a fresh register $r'$ is introduced to count the number of occurrences of $\separator$ that corresponds to the position where the element is extracted. Formally, $\elem^{-1}(\aut)$ is computed as $(\autb \cap \aut_0, true)$ such that $\autb = (R \cup \{r'\}, Q \cup \{pre, post\}, \Sigma', \delta', \{pre\}, \{post\})$, where $\delta'$ comprises the following tuples, 
\begin{itemize}
\item $(pre, a, pre, (\myvec{0}, 0))$ such that $a \in \Sigma$ (thus $a \neq \separator$),  
%
\item $(pre, \separator, pre, (\myvec{0}, 1))$,
%
\item $(pre, \separator, p, (\myvec{v}, 1))$ such that $p \in I$, 
%\item $(pre, \separator, p, (\myvec{0}, 1))$ such that $p \in I$, 
%
\item $(q, a, q', (\myvec{v}, 0))$ such that $(q, a, q', \myvec{v}) \in \delta$ and $a \in \Sigma$ (thus $a \neq \separator$), 
%
\item $(q, \separator, post, (\myvec{0}, 0))$ such that $q \in F$,
%
\item $(post, a, post, (\myvec{0}, 0))$ and $(post, \separator, post, (\myvec{0}, 0))$ where $a \in \Sigma$.
\end{itemize}
Note that in the backward propagation of $\aut$ with respect to an equality $\svarz = \elem(\svarx, it)$, the constraint $r' = it$ is added to the LIA constraint to assert that the value of $r'$ is equal to the index $it$.

%%%%%%%%%%%%%%%%%%%%%%%%%%%%%
\hide{
\begin{example}
Consider the constraints $\svarv = \elem(\svaru, 2) \wedge \svarv = a$. Then the automaton $\aut$ that $\svarv$ belongs to is $q_0\xrightarrow{a}q_1$. The $\elem^{-1}(\aut)$ is $(\autb \cap \aut_0, true)$, where $\autb$ is $pre\xrightarrow[(1)]{\separator}pre\xrightarrow[(0)]{\Sigma}pre\xrightarrow[(1)]{\separator}q_0\xrightarrow[(0)]{a}q_1\xrightarrow[(0)]{\separator}post\xrightarrow[(0)]{\Sigma'}post$. After back propagation, we remove the constraint $\svarv = \elem(\svaru, 1)$ and add the constraints $r' = 2\wedge \svaru \in \autb \cap \aut_0$, where $r'$ is the fresh register introduced in $\autb$. Solving the new constraints, we obtain a solution $\svaru = \separator\separator a \separator$ which represents the sequence $(\epsilon, a)$.
\end{example}
}
%$\myjoin_u$: special case of transducers

\medskip

Finally, we  remark that the operation $\myseqlen(\svarx)$ can be modeled as a CEFA, similarly to the CEFA for $\mystrlen$ in \cite{atva2020}.
%%%%%%%%%%%%%%%%%%%%%%%%%%%%%%%%%%
%%%%%%%%%%%%%%%%%%%%%%%%%%%%%%%%%%
\hide{
\paragraph*{Pre-image of $\myseqlen$.} 
Let $\aut = (R, Q, \Sigma, \delta, I, F)$ be a CEFA. 
The computation of $\myseqlen^{-1}(\aut)$ is similar to that of $\mystrlen^{-1}(\aut)$ in \cite{atva2020}, except that $\myseqlen^{-1}(\aut)$ count the number of occurrences of $\separator$ while $\mystrlen^{-1}(\aut)$ counts the length of the string. Formally, $\myseqlen^{-1}(\aut)$ is computed as $(\autb \cap \aut_0, t_{(p,q)})$ where $t_r = r^{(1)}$ for each $r \in R$, $\autb = (R \cup \{r'\}, \{q\}, \Sigma', \delta', \{q\}, \{q\})$ and $\delta'$ comprises the following tuples,
\begin{itemize}
\item $(q, a, q, (0))$ such that $a \in \Sigma$ (thus $a \neq \separator$),
\item $(q, \separator, q, (1))$.
\end{itemize}
Note that in the backward propagation of $\aut$ with respect to an equality $it = \myseqlen(\svaru)$, the constraint $r' = it$ is added to the LIA constraint to assert that the value of $r'$ is equal to the length of the sequence.

\begin{example}
	Let us consider the constraints $it = \myseqlen(\svaru) - 1\wedge it > 0$. The $\myseqlen^{-1}$ is $(\autb \cap \aut_0, true)$, where $\autb$ is $q\xrightarrow[(1)]{\separator}q\xrightarrow[(0)]{\Sigma}q$. After back propagation, we remove the constraint $it = \myseqlen(\svaru)$ and add the constraints $r' = it\wedge \svaru \in \autb \cap \aut_0$, where $r'$ is the fresh register introduced in $\autb$. Solving the new constraints, we obtain a solution $it = 1$ and $\svaru = \separator\separator$ which represents the sequence $(\epsilon)$. 
\end{example}
}



%% write by denghang ==================================================================
%% write by denghang ==================================================================
\hide{
	\subsection{Special CEFA for String Array}
	%%
	We define a special CEFA for string array. The CEFA uses a preserved character $``\arrseparator"$ as $array$ elements separator (note that $\arrseparator\not\in\Sigma$). The speciality of $\aut_s$ rely on  its construction.
	%%
	\paragraph{Array of fixed length and constant index} Considering an $array$ $A$ with fixed length $m$, we assume that each array element $A[i]$ $(i\in[1,m])$ is constrained by a CEFA $\aut_i\equiv (R_i, Q_i, \Sigma, \delta_i, I_i, F_i, \alpha_i)$. Note that the index of the array starts from 1.  Then we can construct the CEFA of the array $A$ as $\aut_s\equiv (R_s, Q_s, \Sigma\cup\{\arrseparator\}, \delta_s, I_s, F_s, \alpha_s)$, where:
	\begin{itemize}
		% \item $R_s = \bigcup\limits_{1\leq i \leq m} R_i$,
		\item $R_s = \bigcup\limits_{1\leq i \leq m} R_i\cup \myset{r}$ for a fresh register r,
		\item $Q_s = \bigcup\limits_{1\leq i \leq m} Q_i\cup \myset{q_0, q_f}$ for fresh states $q_0$ and $q_f$,
		\item $\delta_s$ is composed of four transitions set
		\begin{itemize}
			\item  $\{(q_0, \arrseparator, q_1, (\myvec{0}, 1))\mid q_1\in I_1\}$,
			\item  $\myset{(q_m, \sharp, q_f, (\myvec{0}, 0))\mid q_m\in F_m}$,
			\item  $\{(q, \arrseparator, q', (\myvec{0}, 1)) \mid q\in F_i, q'\in I_{i+1}, i\in[1,m]\}$ where $ I_{m+1}=\emptyset$,
			\item  and $\{(q, a, q', (\myvec{0},\myvec{v_i},\myvec{0}))\mid (q, a, q', \myvec{v_i})\in \aut_i, i\in[1,m]\}$,
		\end{itemize}
		where $(\myvec{0}, 1)$ is a vector that increases the value of $r$ by 1 and remains other registers unchanged, and $(\myvec{0},\myvec{v_i},\myvec{0})$ is a vector that updates the value of $R_i$ by $\myvec{v_i}$ and remains other registers unchanged. Similarly for $(\myvec{0}, 0)$ and so on.
		\item $I_s = \myset{q_0}$, $F_s = \myset{q_f},$ and $\alpha_s = \bigwedge\limits_{1\leq i \leq m} \alpha_i$.
		% and $\alpha_s = \bigwedge\limits_{1\leq i \leq m} \alpha_i$.
	\end{itemize}
	The resulting CEFA is shown in fig \ref{fig:fixed_len_cons_idx}, where $q_i\in I_i$ and $q_i'\in F_i$ for each $i\in [1,m]$. The main idea of the construction above is using $\arrseparator$ to separate each element of the array and one register to record the index of each element. We introduce two fresh states $q_0$ and $q_f$ to ensure each element is started and ended with at least one $\arrseparator $, which is beneficial to the construction of pre-images of string constraints (see \ref*{sec:pre_image} for more details). The accepting condition $\alpha_s$ is the conjunction of accepting conditions of each element.
	% It is obviously that the accepting word $w$ of $\aut_s$ is in language $\mathcal{L}(\aut_1)\arrseparator\mathcal{L}(\aut_2)\ ... \ \arrseparator \mathcal{L}(\aut_{m}) $.
	\begin{figure}[H]
		\centering
		\usetikzlibrary {automata, positioning, fit}
		\begin{tikzpicture}[initial text =, initial distance=3ex,
			node distance=1.2cm, auto,
			state/.style={circle, draw, inner sep=0pt, minimum size=5mm},
			accepting by double]
			
			\node[state,initial]            (q_0)                {$q_0$};
			\node[state]                    (q_1) [right=of q_0] {$q_1$};
			\node[state]                    (q_1') [right=of q_1] {$q_1'$};
			\node                           (omit)[right=of q_1'] {$\cdots$};
			\node[state]                    (q_m) [right=of omit] {$q_m$};
			\node[state]                    (q_m') [right=of q_m] {$q_m'$};
			\node[state, accepting]         (q_f)  [right=of q_m'] {$q_f$};
			
			\node[draw, fit=(q_1) (q_1'), inner sep=0.5em, label=above:$\aut_1$] [dashed] {};
			\node[draw, fit=(q_m) (q_m'), inner sep=0.5em, label=above:$\aut_m$] [dashed] {};
			
			\path[->]
			(q_0) edge [above] node {$\scriptstyle \arrseparator$} (q_1)
			(q_0) edge [below] node {${\scriptstyle r++}$} (q_1)
			(q_1') edge [above] node {$\scriptstyle \arrseparator$} (omit)
			(q_1') edge [below] node {${\scriptstyle r++}$} (omit)
			(omit) edge [above] node {$\scriptstyle \arrseparator$} (q_m)
			(omit) edge [below] node {${\scriptstyle r++}$} (q_m)
			(q_m') edge [above] node {$\scriptstyle \arrseparator$} (q_f)
			(q_m') edge [below] node {${\scriptstyle r}$} (q_f);
			
		\end{tikzpicture}
		\caption{The CEFA of $A$ with accepting condition $\alpha_s = \bigwedge\limits_{1\leq i \leq m} \alpha_i$}
		\label{fig:fixed_len_cons_idx}
	\end{figure}
	The construction of $\aut_s$ can be used to formalize the array containing constant string. For example, for a 2-length array $A$ constrained by $``a" = \myread(A, 1)\wedge ``b" = \myread(A, 2) \wedge \mylen(A) = 2$ , where $\myread(A, m)$ read the  element of $A$ at constant index $m$ and $\mylen(A)$ is the length of $A$, we can construct the automaton of $A$ as following: \\
	$(\emptyset, \{q_0,q_1,q_1',q_2,q_2'\}, \{a,b,\arrseparator\}, \{q_0\xrightarrow[(1)]{\arrseparator}q_1, q_1\xrightarrow[(0)]{a}q_1', q_1'\xrightarrow[(1)]{\arrseparator}q_2, q_2\xrightarrow[(0)]{b}q_2'\}, \{q_0\}, \{q_2'\}, \top)$. 
	%%
	\paragraph{Array of varible length and index} \todo{TODO}
}%% write by denghang ==================================================================
%% write by denghang ==================================================================



%%%%%%%%%%%%%%%%%%%%%%%%%%%%%%%%%%
%%%%%%%%%%%%%%%%%%%%%%%%%%%%%%%%%%
\hide{

\subsection{Pre-images of String Constraints}\label{sec:pre_image}

In this section, we define the pre-images of string constraint formulae. 
\pagebreak

We let $\aut_{n}$ be the automaton accepting all arrays with length $n$, that is, $\aut_{n} = q_0\xrightarrow{\arrseparator}q_1\cdots\xrightarrow{\arrseparator}q_{n}$ with $q_i\xrightarrow{\Sigma}q_i$ for $0\leq i \leq n-1$. $\aut_*$ is the automaton accepting all arrays with arbitrary length. Then the pre-images of operations are defined as follow.

\paragraph{- $i = \mylen(x)$:} As \cite{atva2020}.

\paragraph{- $x = y\cdot z$:} As \cite{ostrich}.

\paragraph{- $x = \transducer(y)$:} As \cite{ostrich}.

\paragraph{- $i = \mylen(A)$:} The pre-image for the array length function is the CEFA $\aut_{arrlen}=(R, Q, \Sigma\cup\{\arrseparator\}, \delta, I, F, \alpha)$, where
\begin{itemize}
	\item $R = \{r\}$ for a fresh register $r$, $Q = \{q_0, q_f\}$, $\delta = \myset{q_0\xrightarrow[(0)]{\Sigma}q_0, q_0\xrightarrow[(1)]{\sharp}q_0,  q_0\xrightarrow{\sharp}q_1} $, $I = \myset{q_0}$, $F = \myset{q_1}$, and $\alpha \equiv r = i$.
\end{itemize}

\paragraph{- $A = B[m:n]$:} Suppose that $A$ is constrained by $\aut_A=(R_A, Q_A, \Sigma\cup\{\arrseparator\}, \delta_A, I_A, F_A, \alpha_A)$, the pre-image of sub-array operation is defined as $\aut_m\cdot (\aut_A\cap \aut_{m-n})\cdot \aut_{*}$.

\paragraph{- $A = \transducer(B)$:} We consider $A$ and $B$ as string varibles and apply the pre-image computation in \cite{ostrich} directly.

\paragraph{- $x = \myread(A, m)$:} Suppose that $x$ is constrained by the CEFA $\aut_x=(R_x, Q_x, \Sigma, \delta_x, I_x, F_x, \alpha_x)$, the pre-image for the array read function is the CEFA $\aut_{read}=(R, Q, \Sigma\cup\{\arrseparator\}, \delta, I, F, \alpha)$ where
\begin{itemize}
	\item $R = R_x$,
	\item $Q = Q_x\cup\myset{q_0, q_1, \cdots, q_{m}}$,
	\item $\delta$ is the union of the following transitions set:
	\begin{itemize}
		\item $\myset{q_i\xrightarrow[\myvec{0}]{\Sigma}q_i\mid i\in [0, m-1]}$,
		\item $\myset{q_i\xrightarrow[\myvec{0}]{\arrseparator}q_{i+1}\mid i\in[0,m-2]}$,
		\item $\myset{q_{m-1}\xrightarrow[\myvec{0}]{\arrseparator}q_0^x \mid q_0^x\in I_x}$,
		\item $\myset{q\xrightarrow[\myvec{v}]{a}q'\mid q\xrightarrow[\myvec{v}]{a}q'\in \delta_x}$,
		\item $\myset{q_f^x\xrightarrow[\myvec{0}]{\arrseparator}q_m \mid q_f^x\in F_x}$,
		\item $\myset{q_m\xrightarrow[\myvec{0}]{\Sigma\cup\arrseparator} q_m}$,
	\end{itemize}
	\item $I = \myset{q_0}$, $F = \myset{q_m}$, and $\alpha = \alpha_x$.
\end{itemize}

\paragraph{- $A = \mywrite(B, x, m)$:} Suppose that $A$ is constrained by $\aut_A=(R_A, Q_A, \Sigma\cup\{\arrseparator\}, \delta_A, I_A, F_A, \alpha_A)$ and $\aut_A = \aut_{pre}\cdot\aut_{ele}\cdot\aut_{post}$. Then there are $O(|\aut_A|)$ possible concatenation portfolios $(\aut_{pre},\aut_{ele}$, $\aut_{post})$, and for each portfolio satisfying $\aut_{pre}\cap\aut_{m-1}\not=\emptyset$ and $\aut_{ele}\cap\aut_1\not=\emptyset$, we obtain a tuple $((\aut_{pre}\cap\aut_{m-1})\cdot\aut_1\cdot\aut_{post}, \aut_{ele}\cap\aut_1)$ which is one of the pre-image for array write function. The final pre-image is a list of these tuples.

\paragraph{- $x = \myjoin(A,u)$:} Suppose that $x$ is constrained by $\aut_x=(R_x, Q_x, \Sigma, \delta_x, I_x, F_x, \alpha_x)$ and $u=a_1a_2\cdots a_n$, the pre-image for array join function, $\aut_{join}=(R_x, Q_x\cup\myset{q_f}, \Sigma, \delta, I_x, \myset{q_f}, \alpha_x)$, is generated from $\aut_x$ by adding the transition $q\xrightarrow{\arrseparator}q'$ for each $q\xrightarrow{u}q'\in\aut_x$ and $q_f^x\xrightarrow{\arrseparator}q_f$ for each $q_f^x\in F_x$. For $q\xrightarrow{u}q'\in \aut_x$, we mean the state of $\aut_x$ transfer from $q$ to $q'$ after reading the word $u$, that is, there is a path $q\xrightarrow{a_1}q_1\cdots q_{n-1}\xrightarrow{a_n}q'$ in $\aut_x$. 

\paragraph{- $A = \mysplit(x, e)$:} Consider $A$ as string, the semantic of this operation is equal to $A = replaceAll(x, e, \arrseparator)$. We can apply the computation of pre-image of replaceAll function in \cite{ostrich} directly.
}




%%%%%%%%%%%%%%%%%%%%%%%%%%%%%%%%%%%%%%%%%%%%%%%%%%%%%%%%%%%

%\vspace{-2mm}
\section{Implementation and Experiments} \label{sect:exp}

%!TEX root = main.tex

We implement a solver  $\ostrichseq$ on top of the string solver OSTRICH~\cite{CHL+19} and the SMT solver $\princess$~\cite{princess08}, using about 5,000 lines of Scala code. The core functionality, preimage computation for dedicated string operations such as $\mywrite$, $\subseq$, $\elem$ and $\myseqlen$ is implemented in about 2,000 lines of code and is integrated directly into the OSTRICH framework. To support the representation of sequences, we built a lightweight library of automata which treat the separator as a special transition. The library is implemented in about 1,000 lines of code. 
%The remaining code is used for the integration with OSTRICH and for the implementation of the decision procedures for the theory of sequences.

%Moreover, we do experiments to evaluate the effectiveness of our approach.

%The experiment is conducted on a benchmark suite comprising the constraints generated from real-world and hand-crafted templates. We compare our implementation with the state-of-the-art SMT solvers, including Z3, and cvc5. The results show that our approach is efficient and capable of solving constraints that existing solvers cannot.


\subsection{Benchmarks}

To evaluate the effectiveness of $\ostrichseq$, we curate two benchmark suites, $\seqbase$ and $\seqext$,
where $\seqbase$ only uses operations that are directly supported by some existing SMT solvers while $\seqext$ uses some string sequence operations that are not directly supported by any existing SMT solvers.

\smallskip
\noindent\emph{$\seqbase$.}
The $\seqbase$ benchmark suite comprises $\arraylogic$ formulas that contain only the generic sequence operations, that is, sequence operations that can be applied to the elements of any type (not necessarily string), such as sequence length $\myseqlen$, sequence concatenation $\cdot$, subsequence $seqt[it_1, it_2]$, sequence read $\myat(seqt, it)$ and sequence write $seqt[it \rightarrow strt]$.
%
$\seqbase$ is designed to facilitate a fair comparison with existing solvers since they do not support the string-specific sequence operations. %where the sequences and strings are tightly related %(e.g. the constraints that contain $\mysplit_e$ and $\myjoin_u$ that convert between strings and sequences).
$\seqbase$ contains 140 $\arraylogic$ constraints that are generated from three templates. %Specifically,
\begin{itemize}
  \item  \polished{The first template $\svaru_1\in e_1 \wedge \svars_2 = \svars_1[n_1\rightarrow\svaru_1, n_2\rightarrow\svaru_1] \wedge \svaru_2 = \myat(\svars_2, m_1)\cdot\myat(\svars_2, m_2) \wedge \svaru_2 \in e_2 \wedge \myseqlen(\svars_2) < \mystrlen(\svaru_2)$ is utilized to curate 40 random  instances, where
  $\svaru_1$ is a string variable, $\svars_1,\svars_2$ are sequence variables, 
  %$s$ is a random generated sequence constant,
  and $e_1,e_2$ are randomly generated regular expressions.
  In the first 20 instances, $n_1, n_2, m_1, m_2$ are integer variables while in 
  the other 20 instances, they are random integers ranging from 0 to 10.}
  % \hide{Generally speaking, the template writes the same string to different positions of a sequence and then checks the length property and regular properties of elements.}
        %    
  \item \polished{The second template $\svaru_1\in e_1 \wedge \svars_2 = \svars_1[n_1\rightarrow\svaru_1, n_2\rightarrow\svaru_2] \wedge \svars_3 = \svars_2[n_3, len_1]\cdot\svars_2[n_4,len_2] \wedge \svaru_3 = \myat(\svars_3,m_1)\cdot\myat(\svars_3,m_2)\cdot\myat(\svars_3,m_3)\cdot\myat(\svars_3,m_4)\wedge\svaru_3\in e_2 \wedge \myseqlen(\svars_3) < \mystrlen(\svaru_3) + 1$ is utilized to curate 40 random instances, where $\svaru_1,\svaru_2$ are string variables, $\svars_1$, $\svars_2,\svars_3$ are sequence variables, and $e_1,e_2$ are randomly generated regular expressions.
  Similar to the first template,
  $n_i, m_i, len_1,len_2$ for $i\in [1,4]$  are integer variables in the first 20 instances but are random integers ranging from 0 to 10 in the other 20 instances.}
  %As above, we then replace all integer constants in this template with integer variables to produce an additional 20 instances.}
  % \hide{Briefly, the template generates a new sequence from an initial sequence by concatenating two overlapping subsequences and then checks the length and regular properties.}
  \item \polished{The third template $\svars_1 = \svars_0[m_1\rightarrow \svaru_1]\wedge\cdots\wedge\svars_{i} = \svars_{i-1}[m_i\rightarrow \svaru_i]\wedge\svaru_1=\myat(\svars_1, n_1)\wedge\cdots\wedge\svaru_j=\myat(\svars_j, n_j)\wedge\svaru_1\in e_1 \wedge\cdots\wedge \svaru_j\in e_j$ is utilized to curate 60 random instances, where $i$ and $j$ are randomly selected from the range $0$ to $20$, $\svars_0 \cdots \svars_i$ are sequence variables,  $\svaru_1\cdots \svaru_j$ are string variables, $e_1\cdots e_j$ are randomly generated regular expressions, and 
  $m_1,\cdots m_i$, $n_1\cdots n_j$ are randomly generated integers constant ranging from 0 to 5. In this template, write ($\svars[\cdot\rightarrow\cdot]$) and read ($\myat(\svars,\cdot)$) operations are performed randomly on a string sequence, and the strings read from the sequence are checked against some regular expressions.}
  % \item The third template $\svars_1 = \svars[m_1\rightarrow \svaru_1]\wedge\svars_2 = \svars_1[m_2\rightarrow \svaru_2]\wedge\cdots\wedge\svaru_1=\myat(\svars_1, n_1)\wedge\svaru_2=\myat(\svars_2, n_2)\wedge\cdots\wedge\svaru_1\in e_1 \wedge \svaru_2\in e_2 \wedge \cdots$ is utilized to generate 60 random instances, where $\svaru_1$, $\svaru_2$ are string variables, $\svars$, $\svars_1$, $\svars_2$ are sequence variables,  $e_1$, $e_2$ are randomly generated regular expressions, and 
  % $m_1$, $m_2$, $n_1$, and $n_2$ are randomly generated  integers constant. In this template, write ($s[\cdot\rightarrow\cdot]$) and read($\myat$) operations are performed randomly on a string sequence, and the strings read from the sequence are checked against some regular expressions.
\end{itemize}



%The benchmark contains 100 constraints generated from 3 different templates. The templates are aimed to cover the basic operations of string sequences. The templates and corresponding number of instances are as follows:

%Note that the regular properties of the sequence are defined by concatenating the first and the second string element in the sequence and checking whether the concatenated string satisfies a regular expression. This benchmark is designed to cover the basic operations of string sequences. 

\paragraph*{$\seqext$.}
The $\seqext$ benchmark suite comprises 60 $\arraylogic$ formulas that use specific string sequence operations, namely, $\myfilter_e$, $\mymatchall_e$, $\mysplit_e$ and $\myjoin_u$. Among them, 19 formulas are manually generated for the unit tests of $\ostrichseq$, and 41 formulas are generated from the real-world JavaScript programs collected from GitHub. Moreover, to compare with SMT solvers that do not support the considered string sequence operations, we rewrite them using the basic string operations whenever possible. ($\mymatchall_e$ is not rewritten since it cannot be expressed using existing string operations.) 
%We partition these formulas into two categories: $\withoutmatch$ and $\withmatch$. The former contains constraints that do not use $\mymatchall_e$, while %the $\withmatch$ category 
%the latter uses it. %contains the constraints that use the operation $\mymatchall_e$.  

%60 constraints generated by hand. 
%41 of them are transformed from the real-world JavaScript programs snippets, and 19 of them are collected from the unit tests of our solver. 
%The benchmark is designed to cover the extended operations of string sequences. 

%%%%%%%%%%%%%%%%%%%%%%%%%%%%%%%%%%%%%
\hide{
  Since we support some extended functionalities of string sequences such as \verb|str.split| and \verb|seq.join|, we divide the benchmark suite into two parts: \textbf{Seq-Base} and \textbf{Seq-Ext}. The \textbf{Seq-Base} benchmark contains the constraints with only the basic operations of string sequences, such as \textbf{length, concatenation, subsequence, read and write}. Note that we also use write in the benchmark, which can be transformed to concatenation, subsequence and length.
  The \textbf{Seq-Ext} benchmark contains the constraints with the extended operations of string sequences, such as \textbf{split, join, match and filter}.

  \begin{itemize}
    \item \textbf{Seq-Base}: The benchmark contains 100 constraints generated from 3 different templates. The templates are aimed to cover the basic operations of string sequences. The templates and corresponding number of instances are as follows:
          \begin{itemize}
            \item \texttt{Template 1}: 20 instances, which write the same string to different positions of the sequence and then check the regular properties of the sequence.

            \item \texttt{Template 2}: 20 instances, which generate a new sequence from the initial sequence by concatenating two overlapping subsequences and then checking the regular properties of the sequence.
            \item \texttt{Template 3}: 60 instances, which randomly read and write to a sequence of strings and then check the regular properties of the read strings.
          \end{itemize}
          Note that the regular properties of the sequence are defined by concatenating the first and the second string elements in the sequence and checking whether the concatenated string satisfies a regular expression. This benchmark is designed to cover the basic operations of string sequences.
    \item \textbf{Seq-Ext}: The benchmark contains 60 constraints generated by hand. 41 of them are transformed from real-world JavaScript program snippets, and 19 of them are collected from the unit tests of our solver. The benchmark is designed to cover the extended operations of string sequences.
  \end{itemize}
}
%%%%%%%%%%%%%%%%%%%%%%%%%%%%%%%%%%%%%

\subsection{Experimental results}
%We compare with cvc5, Z3, Z3-noodler, $\ostrich$ and $\princessarr$ on them.

We conduct experiments by comparing $\ostrichseq$ with SOTA solvers, including cvc5, Z3, Z3-noodler, $\ostrich$ and $\princessarr$ (referring to an array-theory-based solver implemented within $\princess$). $\seco$~\cite{JezLMR23} is excluded %in the comparison 
due to soundness issues. Additionally, \verb|foldleft| and \verb|map| operations are not used to simulate $\myjoin$, as most constraints return unknown under such usage in Z3.
%
On $\seqbase$, we evaluate $\ostrichseq$ along with sequence solvers cvc5, Z3 and $\princessarr$. For $\seqext$, we rewrite the constraints using only the basic string operations and compare $\ostrichseq$ with string solvers cvc5, Z3, Z3-noodler and $\ostrich$.

All experiments are run on a 24 $\times$ 3.1 GHz-core server with 185 GB RAM, where each solver is started as a single thread, with a 60 seconds time limit and without the memory limit.

\begin{table}[t]
\setlength{\tabcolsep}{6pt}
%\renewcommand{\arraystretch}{1.1}
%\vspace{-6mm}
  \centering
    \caption{\polished{Experimental results on $\seqbase$.}
  	%$\ostrichseq$ can solve all the constraints in $\seqbase$ and is more efficient than cvc5 and Z3. $\princessarr$ is slightly more efficient than $\ostrichseq$.
  } 
    \begin{tabular}{| c | c | c | c | c |}
      \hline
                      & cvc5        & Z3    & $\princessarr$ & $\ostrichseq$ \\
      \hline
      sat             & 48          & 14    & \textbf{52}    & \textbf{52}   \\
      \hline
      unsat           & 87          & 59    & 77             & \textbf{88}   \\
      \hline
      solved          & 135         & 73    & 129            & \textbf{140}  \\
      \hline
      unknown/timeout & 5           & 67    & 11             & \textbf{0}    \\
      \hline
      avg. time (s)   & 3.53        & 24.36 & 6.84           & \textbf{3.23}          \\
      \hline
    \end{tabular} %
%  \end{center}
  \label{tab:seq-base}
%\vspace{-8mm}
\end{table}


The results on $\seqbase$ are reported in Table~\ref{tab:seq-base}.
We can see that $\ostrichseq$ outperforms all the other solvers on $\seqbase$ in terms of both the number of solved instances and the solving time per instance. Specifically,   $\ostrichseq$ solves all the \polished{140} instances in $\seqbase$, while cvc5, Z3 and $\princessarr$ only solve \polished{135}, \polished{73} and \polished{129} instances, respectively. Moreover,  $\ostrichseq$ is more efficient than cvc5, Z3 and $\princessarr$ ($\ostrichseq$ spends \polished{3.23} seconds in solving each instance from $\seqbase$ on average, while cvc5, Z3 and $\princessarr$ spend \polished{3.25}, \polished{24.36}, and \polished{6.84} seconds, respectively). 

% \hide{However, $\ostrichseq$ is slightly slower than $\princessarr$, which solves all the instances in $\seqbase$ in 4.04 seconds on average. This is because $\princessarr$ utilizes the optimization techniques that are specialized to the array theory, making it more efficient for solving constraints that involve only read and write operations.}


The results on $\seqext$ are reported in Table~\ref{tab:seq-ext}. Since the SOTA sequence solvers do not support $\myfilter_e$, $\mymatchall_e$ or $\mysplit_e$ directly, the $\arraylogic$ constraints are rewritten to string constraints (instead of solving the $\arraylogic$ constraints directly) when we compare with them on $\seqext$.
% we fulfill the comparison with string solvers on $\seqext$ by 
%call string solvers use them to solve the string constraints generated from the $\arraylogic$ constraints (instead of solving the $\arraylogic$ constraints directly). 
Moreover, while we can use {\sf str.replace\_re\_all} to encode $\myfilter_e$ and $\mysplit_e$, the encoding of $\mymatchall_e$ requires finite-state transducers which have not been supported by most SOTA string solvers. %(except for {\sf OSTRICH}, \polished{but we do not implement this transformation in $\ostrich$, as it has already been applied in $\ostrichseq$}). 
%\denghang{{\sf OSTRICH} offer built-in support for finite-state transducers, but defining a transducer within an smtlib2.7 file is not feasible}). 
As a result, we do \emph{not} generate the string constraints for the 12 instances in $\seqext$ that contain $\mymatchall_e$, which are denoted by 
%We use 
$\seqext\mbox{\sf -M}$, whereas the rest 48 instances are %to denote the 12 instances in $\seqext$ that contain $\mymatchall_e$ and 
denoted by $\seqext\mbox{\sf -N}$. %to denote the other 48 instances in $\seqext$.


\begin{table}[t]
\setlength{\tabcolsep}{2pt}
%\renewcommand{\arraystretch}{1}
%\vspace{-6mm}
\caption{Experimental results on $\seqext$.
	% $\withmatch$ refers to the constraints that contain the operation $\mymatchall_e$, while $\withoutmatch$ refers to the constraints that do not contain the operation $\mymatchall_e$. The results show that $\ostrichseq$ can solve all the constraints in $\seqext$ and is more efficient than cvc5 and $\ostrich$. Z3 and Z3-noodler solve only a few constraints but are much faster than others.
}
%\vspace{-3mm}
\centering
    \begin{tabular}{| c| c| c | c | c | c| c|}
      \hline
                                              &                 & cvc5  & Z3   & Z3-noodler    & $\ostrich$ & $\ostrichseq$ \\
      \hline
      \multirow{3}{*}{$\seqext\mbox{\sf -N}$} & sat             & 19    & 4    & 4             & 20         & \textbf{23}   \\
      \cline{2-7}
                                              & unsat           & 21    & 6    & 4             & 18         & \textbf{25}   \\
      \cline{2-7}
                                              & solved          & 40    & 10   & 8             & 38         & \textbf{48}   \\
      \cline{2-7}
                                              & unknown/timeout & 8     & 38   & 40            & 10         & \textbf{0}    \\
      \cline{2-7}
                                              & avg. time (s)   & 10.83 & 0.13 & \textbf{0.05} & 15.92      & 3.53          \\
      \hline
      \multirow{3}{*}{$\seqext\mbox{\sf -M}$} & sat             & -     & -    & -             & -          & \textbf{12}   \\
      \cline{2-7}
                                              & unsat           & -     & -    & -             & -          & \textbf{0}    \\
      \cline{2-7}
                                              & solved          & -     & -    & -             & -          & \textbf{0}    \\
      \cline{2-7}
                                              & unknown/timeout & -     & -    & -             & -          & \textbf{0}    \\
      \cline{2-7}
                                              & avg. time (s)   & -     & -    & -             & -          & \textbf{1.77} \\
      \hline
    \end{tabular}
  \label{tab:seq-ext}
%\vspace{-8mm}
\end{table}


From Table~\ref{tab:seq-ext}, we can see that Z3 and Z3-noodler solve only 10 and 8 out of 48 constraints in $\seqext\mbox{\sf -N}$, respectively, and cvc5 and $\ostrich$ solve 40 and 38 instances, respectively,
indicating that the complexity of the string constraints generated from the $\arraylogic$ constraints. %are too complex to be solved by string solvers successfully in general. 
On the other hand, $\ostrichseq$ solves all of them, demonstrating its superiority in solving the $\arraylogic$ constraints with string sequence operations. For efficiency, %the solving time per instance, 
$\ostrichseq$ spends 3.53 seconds per instance on average, significantly less than cvc5 and $\ostrich$ (10.83 and 15.92 seconds on average, respectively). For the constraints in $\seqext\mbox{\sf -M}$, $\ostrichseq$ successfully solves all 12 instances with an average time of 1.77 seconds.
As aforementioned, for the constraints in $\seqext\mbox{\sf -M}$, we do not generate string constraints or compare against cvc5, Z3, Z3-noodler and $\ostrich$, since these constraints cannot be expressed using only the basic operations of string theory.
%are unable to solve any of these instances, as $\mymatchall_e$ cannot be rewritten using only the basic operations of string theory.
%In contrast, cvc5, Z3, Z3-noodler are unable to solve any of these instances, as $\mymatchall_e$ cannot be rewritten using only the basic operations of string theory.



%The results of the Seq-Base benchmark are shown in Table~\ref{tab:seq-base}. The results show that our solver is slower than Z3 and cvc5 in some cases, but it is still able to solve all the constraints. 

%%%%%%%%%%%%%%%%%%%%%%%%%%%%%%%%%%%%
%%%%%%%%%%%%%%%%%%%%%%%%%%%%%%%%%%%%
% \hide{
%   The results of the Seq-Ext benchmark are shown in Table~\ref{tab:seq-ext}. Since the extended operations are not supported by Z3 and cvc5, we rewrite the constraints to the basic operations of string theory, where we use string $'\#aa\#bb'$ to stand for the sequence $['aa', 'bb']$. \textbf{filter} and \textbf{match} are not rewrote since they are transducer which is not supported even in string theory. The results show that our solver OSTRICH outperforms Z3 and cvc5 in this benchmark, solving all the constraints with a competetive time. 28 unknown constraints of cvc5 and Z3 are due to filter and match. 17 unknown constraints of cvc5 and 25 unknown constraints of Z3 are due to the axioms we introduced to simulate the extended operations such as \textbf{split} and \textbf{join}.
% }
%%%%%%%%%%%%%%%%%%%%%%%%%%%%%%%%%%%%
%%%%%%%%%%%%%%%%%%%%%%%%%%%%%%%%%%%%




%%%%%%%%%%%%%%%%%%%%%%%%%%%%%%%%%%%%%%%%%%%%%%%%%%%%%%%%%%%

%\vspace{-2mm}
\section{Related work}\label{sect:related}

%!TEX root = main.tex


%\smallskip
\noindent
{\bf String constraint solving.} 
String constraint solving has received considerable attention in the literature. Amadini \cite{Amadini23} provides a comprehensive survey of the literature up to 2021. In particular, %extensive documentation of this diversity, and classifies
approaches to string constraint solving are classified into three main categories: automata-based (relying on finite state automata to represent the domain of string variables
and to handle string operations, e.g., \cite{Abdulla14,HJLRV18,CCH+18,Vijay-length,LB16,BerzishKMMDNG20,ChenCHHLS24}), word-based (algebraic approaches based on systems of word equations, e.g., \cite{cvc4,BarbosaBBKLMMMN22,Z3,Z3str2,Z3str3,Z3str4,TCJ16,S3,BerzishKMMDNG20,ChenCHHLS24} 
), 
and unfolding-based approaches (explicitly reducing each string into a number of contiguous elements denoting its
characters, e.g., \cite{SFPS17,KGAGHE13,SAHMMS10,LG13,AGS18,DEKMNP19}). 

%Generally speaking, 
From another perspective, there are two lines of work aiming for building practical string solvers: %could essentially be classified into two categories. 
(1) one could support as many constraints as possible, but primarily resort to heuristics, offering no completeness/termination guarantee. 
%This is a realistic approach since, as mentioned above, the problem involving both string and integer data types is in general undecidable. 
Many solvers %for string constraints 
are usually implemented as string
theory solvers inside SMT solvers, %, such as cvc5 [9] or Z3 [31], 
allowing combination with other theories, most commonly the theory of integers for string lengths. Some (non-exhaustive) examples include CVC4/5 \cite{cvc4,BarbosaBBKLMMMN22}, Z3 \cite{Z3,BTV09}, Z3-str2/3/4 \cite{Z3str2,Z3str3,Z3str4}, S3(P)~\cite{TCJ16,S3}, Trau~\cite{Abdulla17} (or its variants Trau+ \cite{AbdullaA+19}), %and
%Z3-Trau [9]), %ABC [10], 
Slent~\cite{WC+18}, etc.  
%
%
%i) Automata-based approaches: .
%(ii) Word-based approaches: .
%(iii) Unfolding-based approaches: 
%String constraint solving has received considerable attentions in literature. Generally speaking, there are two lines of work aiming for building 
%practical string solvers: %could essentially be classified into two categories. 
%(1) one strives to support as many constraints as possible, but primarily resort to heuristics, offering no completeness/termination guarantee. 
More recent ones include Z3str3RE \cite{BerzishKMMDNG20} %Z3-Trau [1], Z3str4 [30],  
%This is a realistic approach since, as mentioned above, the problem involving both string and integer data types is in general undecidable. 
% 
%Completeness guarantees are, however, valuable since the performance of heuristics can be difficult to predict. 
and Z3-Noodler \cite{ChenCHHLS24} %is a fork of Z3 that replaces its string theory solver with a custom solver implementing the recently introduced 
(which is based on stabilization-based algorithm \cite{BlahoudekCCHHLS23,ChenCHHLS23} for solving word equations with regular constraints). 
%
%An extensive experimental evaluation shows that Z3-Noodler is a fully-fledged solver that can compete with state-of-the-art solvers, surpassing them by far on many benchmarks. Moreover, it is often complementary to other solvers, making it a suitable choice as a candidate to a solver portfolio.
%
(2) The second approach is to develop solvers for decidable fragments. 
%supporting both strings and integers %
Solvers in this category include Norn \cite{Abdulla14}, SLOTH \cite{HJLRV18} and
OSTRICH \cite{CCH+18,CHL+19}, which are usually based on complete decision procedures (e.g. \cite{Vijay-length,AbdullaA+19,LB16}) 

%The fragment without complex string operations (e.g. replaceAll and
%finite transducers, but length) can be handled quite well by Norn. The fragment without
%length constraints (but replaceAll and finite transducers) can be handled effectively by
%OSTRICH and SLOTH. Moreover, most existing solvers that belong to the first category do not support complex string operations like replaceAll and finite transducers
%as well. This motivates the following problem: provide a decision procedure that supports both string and integer data type, with completeness guarantee and meanwhile
%admitting efficient implementation

We also mention that there are solvers which emphasize certain aspects of string constraint solving by providing dedicated methods or optimization. These include, for instance, those for more expressive regular expressions \cite{BerzishDGKMMN23,ChenFHHHKLRW22},  regex-counting~\cite{HuW23},  %A decision procedure for string constraints with 
string/integer conversion and flat regular constraints~\cite{WuCWXZ24}, Not-Substring Constraint~\cite{AbdullaACDHHTWY21}, integer data type~\cite{atva2020}, etc. 

%	Parosh Aziz Abdulla, Mohamed Faouzi Atig, Yu-Fang Chen, Bui Phi Diep, Lukás Holík, Denghang Hu, Wei-Lun Tsai, Zhilin Wu, Di-De Yen:Solving Not-Substring Constraint withFlat Abstraction. APLAS 2021: 305-320
%\tl{some other work popl22, zhilin's work with Yufang, I cannot access dblp.}

\smallskip
\noindent
{\bf Sequences as an extension of strings.} Strings and sequences are closely related, but may exhibit in different forms. For instance, a string can be considered as a sequence over a finite alphabet. 
Jez et al.~\cite{JezLMR23} recently studied sequence theories which are an extension of theories of strings with an infinite alphabet of letters, together with a corresponding alphabet theory (e.g. linear integer arithmetic). They provide decision procedures based on parametric automata/transducers, giving rise to a sequence constraint solver $\seco$. Our work is different in that we essentially work on a sequence theory where the element is instantiated by the string type. 

%
%A. Jez etc.~\cite{JezLMR23} present a formal and practical study of decision procedures for sequence theories logical theories that generalize string reasoning to sequences over infinite alphabets, such as integers or strings themselves. The main idea is to introduce parametric automata (PA) and parametric transducers as generalizations of symbolic automata, enabling expressive yet decidable formulations of sequence constraints. To support practical verification tasks, they develop $\seco$, a sequence constraint solver based on the parametric automata. Unfortunately, despite the power of parametric automata, the implementation of $\seco$ is not yet mature enough to be used in practice, which has some soundness issues on our benchmarks.

%The CVC5 team propose using an SMT theory of sequences to represent vectors more naturally and flexibly~\cite{cvc5_seq}. They develop two proof calculi for this theory. (1) A base calculus, adapted from existing techniques for strings, focusing on operations like sequence concatenation, extraction, nth, and update, (2) an extended calculus, which incorporates array-style reasoning (e.g., read-over-write axioms) directly into the sequence theory, enabling more efficient handling of nth and update.
%The experiment result shows that the extended calculus significantly outperforms the base calculus. 

%Z3~\cite{Z3} is a well-known SMT solver that has supported sequence theory but has not published a literature on it. The documentation of Z3 \cite{Z3_seq} provides a brief overview of its sequence theory, which support all the operations used in CVC5 except for the update operation. But it supports two more general operations: map and foldleft.

\smallskip
\noindent
{\bf Theory of arrays.}
There are many studies on solving formulas in the theory of arrays or its extensions, for instance, \cite{SBDL01,BMS06,HIV08,FalkeMS13,DacaHK16,FarinierDBL18}, which are also related to the sequence theory. Note that in these studies, the elements are considered to be either of a generic type or of integer type. In other words, the theory of arrays where the elements are of the string type has not been defined and investigated therein. 
%The theory of arrays is a fundamental component in software verification, particularly for modeling memory and data structures. 
%For instance, \cite{FarinierDBL18} proposes a new approach to array formula simplification based on a new dedicated data structure;  \cite{DacaHK16} proposed array fold logic (AFL), which is an extension to the quantifier-free theory of integer arrays to express counting. The satisfiability of AFL is shown by a reduction to counter machines.

%The properties expressible in Array Folds Logic (AFL) include statements such as "the first array cell contains the array length," and "the array contains equally many minimal and maximal elements." These properties cannot be expressed in quantified fragments of the theory of arrays, nor in the theory of concatenation. Using reduction to counter machines, we show that the satisfiability problem of AFL is PSPACE-complete, and with a natural restriction the complexity decreases to NP. We also show that adding either universal quantifiers or concatenation leads to undecidability.
%AFL contains terms that fold a function over an array. We demonstrate that folding, a well-known concept from functional languages, allows us to concisely summarize loops that count over arrays, which occurs frequently in real-life programs. We provide a tool that can discharge proof obligations in AFL, and we demonstrate on practical examples that our decision procedure can solve a broad range of problems in symbolic testing and program verification.
% 


%\cite{FarinierDBL18}
%The theory of arrays has a central place in software verification due to its ability to
%model memory or data structures. Yet, this theory is known to be hard to solve in both
%theory and practice, especially in the case of very long formulas coming from unrolling-based
%verification methods. Standard simplification techniques à la read-over-write suffer from
%two main drawbacks: they do not scale on very long sequences of stores and they miss
%many simplification opportunities because of a crude syntactic (dis-)equality reasoning. We
%propose a new approach to array formula simplification based on a new dedicated data
%structure together with original simplifications and low-cost reasoning. The technique is
%efficient, scalable and it yields significant simplification. The impact on formula resolution
%is always positive, and it can be dramatic on some specific classes of problems of interest,
%e.g. very long formula or binary-level symbolic execution. While currently implemented as
%a preprocessing, the approach would benefit from a deeper integration in an array solver
%
%Decision procedures for the theory of arrays. Surprisingly, there have been relatively few
%works on the efficient handling of the (basic) theory of arrays. Standard symbolic approaches
%for pure arrays complement symbolic read-over-write preprocessing [22, 5, 2] with enumeration
%on (dis-)equalities, yielding a potentially huge search space. New array lemmas can be added
%on-demand or incrementally discovered through an abstraction-refinement scheme [9]. Another
%possibility is to reduce the theory of arrays to the theory of equality by systematic “inlining” of
%the array axioms to remove all store operators, at the price of introducing many case-splits The
%encoding can be eager [24] or lazy [9]. Our method generalizes previous preprocessing [22, 2]
%and is complementary to complete resolution methods [9, 24]. Note also that our approach
%could benefit from being integrated directly within such a complete resolution method, allowing
%incremental simplification all along the resolution process.
%Decision procedures have also been developed for expressive extensions of the array theory,
%such as arrays with extensionality (i.e. equality over whole arrays) or the array property fragment
%[7], which enables limited forms of quantification over indexes and arithmetic constraints. These
%extensions aim at increasing expressiveness and they do not focus so much on practical efficiency.
%Our method can also be applied to these settings (as row are still a crucial issue), even though
%it will not cover all difficulties of these extensions.

%. 




%Recently,many tools for solving string constraints have been developed, motivated
%mainlybytechniquesforfindingsecurityvulnerabilitiessuchasSQLinjectionorcross-
%site scripting(XSS)inwebapplications[34,35,36].
%
%String solving has also found its
%applications in,e.g.,analysis of access user policies in Amazon Web Services [26,8,39]
%or smart contracts [7].






%Sequence theories are an extension of theories of strings with an infinite alphabet of letters, together with a corresponding alphabet theory (e.g. linear integer arithmetic). Sequences are natural abstractions of extendable arrays, which permit a wealth of operations including append, map, split, and concatenation. In spite of the growing amount of tool support for theories of sequences by leading SMT-solvers, little is known about the decidability of sequence theories, which is in stark contrast to the state of the theories of strings. We show that the decidable theory of strings with concatenation and regular constraints can be extended to the world of sequences over an alphabet theory that forms a Boolean algebra, while preserving decidability. In particular, decidability holds when regular constraints are interpreted as parametric automata (which extend both symbolic automata and variable automata), but fails when interpreted as register automata (even over the alphabet theory of equality). When length constraints are added, the problem is Turing-equivalent to word equations with length (and regular) constraints. Similar investigations are conducted in the presence of symbolic transducers, which naturally model sequence functions like map, split, filter, etc. We have developed a new sequence solver, SeCo, based on parametric automata, and show its efficacy on two classes of benchmarks: (i) invariant checking on array-manipulating programs and parameterized systems, and (ii) benchmarks on symbolic register automata.

%Sequences are an extension of strings, wherein elements might range over an infinite domain (e.g., integers, strings, and even sequences themselves). Sequences are ubiquitous and commonly used data types in modern programming languages. They come under different names, e.g., Python/Haskell/Prolog lists, Java ArrayList (and to some extent Streams) and JavaScript arrays. Crucially, sequences are extendable, and a plethora of operations (including append, map, split, filter, concatenation, etc.) can naturally be defined and are supported by built-in library functions in most modern programming languages.
%
%Various techniques in software model checking [30] — including symbolic execution, invariant generation — require an appropriate SMT theory, to which verification conditions could be discharged. In the case of programs operating on sequences, we would consequently require an SMT theory of sequences, for which leading SMT solvers like Z3 [6, 38] and cvc5 [4] already provide some basic support for over a decade. The basic design of sequence theories, as done in Z3 and cvc5, as well as in other formalisms like symbolic automata [15], is in fact quite natural. That is, sequence theories can be thought of as extensions of theories of strings with an infinite alphabet of letters, together with a corresponding alphabet theory, e.g. Linear Integer Arithmetic (LIA) for reasoning about sequences of integers. Despite this, very little is known about what is decidable over theories of sequences.

%%%%%%%%%%%%%%%%%%%%%%%%%%%%%%%%%%%%%%%%%%%%%%%%%%%%%%%%%%%

\section{Conclusion} \label{sect:conc}

In this paper, we have proposed $\arraylogic$, a logic of string sequences, which supports a wealth of string and sequence operations---especially their interactions---commonly in real-word string-manipulation programs. We have provided a decision procedure for the straight-line fragment, which is implemented as a new tool $\ostrichseq$. The experiments on both hand-crafted and real-world benchmarks demonstrate that $\ostrichseq$ is an effective and promising tool for testing, analysis and verification of string-manipulation programs in practice. 

%
%and study its satisfiability. We show that, while it is undecidable in general,  the decidability can be recovered by restricting to the straight-line fragment. This is shown by encoding a string sequence into a string and accordingly, encoding string sequence operations into corresponding string operations. %providing associated 
%We provide pre-image computation for the resulting string operations with respect to automata, effectively casting it into the generic OSTRICH string constraint solving framework. %to which existing decision procedure for strings and integer constraints can be applied.  
%%
%We implement the new decision procedure as a tool $\ostrichseq$,  and carry out experiments on benchmark constraints generated from both real-world and hand-crafted JavaScript programs. The experiments confirm the efficacy of our approach.
\newpage

\bibliographystyle{abbrv}
\bibliography{ref}

\end{document}